\documentclass{amsart}
%\documentclass{amsbook}


%\documentstyle[amsmath,amssymb,amscd,graphicx,psfrag,epsf,amsthm]{article}



%\documentclass[10pt,leqno,amssymbols]{amsbook}

%\documentstyle[12pt,twoside]{amsart}
%\documentstyle[amstex,12pt,leqno]{book}
%\documentstyle[amstex,12pt,leqno]{amsart}
%\documentstyle[12pt,leqno]{amsbook}


\usepackage{amssymb}
\usepackage[all]{xy}
\usepackage{hyperref}


\usepackage{enumitem}   %from SK, fixes \enumerate
% \setlist[enumerate]{itemsep=.2em,topsep=.2em,leftmargin=2.3em,
% label={(\thethm.\arabic*)}}
\renewcommand{\labelenumi}{{\rm(\thethm.\arabic{enumi})}}
\renewcommand{\labelenumii}{{\rm(\thethm.\arabic{enumi}.\alph{enumii})}}
\setlist[enumerate]{itemsep=.2em,topsep=.2em,leftmargin=1.25em,itemindent=2.0em}

%%%%%%%%%%%%%%%%%%%%%%%%%%%%%%%%%%%%%%%%%%%%%%%%%%%%%%%%%%%%%%%%%%%%%%

% Definitions for new environments
% theorem style plain --- default
%\newtheorem{thm}{Theorem}[chapter]
\newtheorem{thm}{Theorem}%[section]
\newtheorem{mainthm}[thm]{Main Theorem}
\newtheorem{question}[thm]{Question}
\newtheorem{lem}[thm]{Lemma}
\newtheorem{cor}[thm]{Corollary}
\newtheorem{cors}[thm]{Corollaries}
\newtheorem{prop}[thm]{Proposition}
\newtheorem{crit}[thm]{Criterion}
\newtheorem{conj}[thm]{Conjecture}
\newtheorem{principle}[thm]{Principle} %!!!!!!!!!!!!!!!!!!!!!!
\newtheorem{complem}[thm]{Complement}%!!!!!!!!!!!!!!!!!!!!!!
\newtheorem{pitfall}[thm]{Pitfall}%!!!!!!!!!!!!!!!!!!!!!!




\theoremstyle{definition}
\newtheorem{defn}[thm]{Definition}
\newtheorem{condition}[thm]{Condition}
\newtheorem{say}[thm]{}
\newtheorem{exmp}[thm]{Example}
\newtheorem{exmps}[thm]{Examples}
\newtheorem{hint}[thm]{Hint}
\newtheorem{exrc}[thm]{Exercise}
\newtheorem{prob}[thm]{Problem}
\newtheorem{const}[thm]{Construction}   %!!!!!!!!!!!!!!!!
\newtheorem{ques}[thm]{Question}    %!!!!!!!!!!!!!!!!!!!!
\newtheorem{alg}[thm]{Algorithm}
\newtheorem{rem}[thm]{Remark}
\newtheorem{rems}[thm]{Remarks}  
\newtheorem{aside}[thm]{Aside}          
%\renewcommand{\theremark}{} 
\newtheorem{note}[thm]{Note}            %\renewcommand{\thenote}{} 
\newtheorem{summ}[thm]{Summary}         %\renewcommand{\thesumm}{} 
\newtheorem*{ack}{Acknowledgments}      % \renewcommand{\theack}{} 
\newtheorem{notation}[thm]{Notation}   
\newtheorem{warning}[thm]{Warning}  
\newtheorem{defn-thm}[thm]{Definition--Theorem}  %!!!!!!!!!!!!!!!!!!!!!!!!
\newtheorem{defn-lem}[thm]{Definition--Lemma}  %!!!!!!!!!!!!!!!!!!!!!!!!
\newtheorem{convention}[thm]{Convention}  %!!!!!!!!!!!!!!!!!!!!!!!!!!!
\newtheorem{keyidea}[thm]{Key idea} 
\newtheorem{expl}[thm]{Explanation}   
\newtheorem{meth}[thm]{Method}   
\newtheorem{assumption}[thm]{Assumption}   
\newtheorem{assertion}[thm]{Assertion}   
\newtheorem{comments}[thm]{Comments}   
\newtheorem{task}[thm]{Task}

\theoremstyle{remark}
\newtheorem{claim}[thm]{Claim}  
\newtheorem{case}{Case} 
\newtheorem{subcase}{Subcase}   
\newtheorem{step}{Step}
\newtheorem{approach}{Approach}
%\newtheorem{principle}{Principle}
\newtheorem{fact}{Fact}
\newtheorem{subsay}{}

%%%%%%%%%%%%%%%%%%%%%%%%%%%%%%%%%%%%%%%%%%%%%%%%%%%%%%%%%%%%%%%%%%%%%%
\setcounter{section}{0}
 %\documentstyle[12pt]{book}
%\setlength{\topmargin}{0in}
%\setlength{\headsep}{0in}
%\setlength{\headheight}{0in}
%\setlength{\textheight}{8.25in}
%\setlength{\parindent}{0in}
%\pagestyle{empty}


\let \cedilla =\c
\renewcommand{\c}[0]{{\mathbb C}}  
\let \crossedo =\o
\renewcommand{\o}[0]{{\mathcal O}} 
%\renewcommand{\o}[0]{{\mathcal O}} 
\newcommand{\z}[0]{{\mathbb Z}}
\newcommand{\n}[0]{{\mathbb N}}
\let \ringaccent=\r  %%% \r  shorthand for 'ring accent'
\renewcommand{\r}[0]{{\mathbb R}} 

\renewcommand{\a}[0]{\hat{\mathbb A}} 

\newcommand{\dd}[0]{{\mathbb D}}
\newcommand{\s}[0]{{\mathbb S}}
\newcommand{\h}[0]{{\mathbb H}}
\newcommand{\p}[0]{{\mathbb P}}
\newcommand{\f}[0]{{\mathbb F}}
\newcommand{\q}[0]{{\mathbb Q}}
\newcommand{\map}[0]{\dasharrow}
\newcommand{\qtq}[1]{\quad\mbox{#1}\quad}
\newcommand{\spec}[0]{\operatorname{Spec}}
\newcommand{\pic}[0]{\operatorname{Pic}}
\newcommand{\pics}[0]{\operatorname{\mathbf{Pic}}}
\newcommand{\picl}[0]{\operatorname{Pic}^{\rm loc}}
\newcommand{\picel}[0]{\operatorname{Pic}^{\rm et-loc}}
\newcommand{\pical}[0]{\operatorname{Pic}^{\rm an-loc}}
\newcommand{\pica}[0]{\operatorname{Pic}^{\rm an}}
\newcommand{\picls}[0]{\operatorname{\mathbf{Pic}}^{\rm loc}}
\newcommand{\piclf}[0]{\operatorname{{\it Pic}}^{\rm loc}}
\newcommand{\piclo}[0]{\operatorname{\mathbf{Pic}}^{\rm loc-\circ}}
\newcommand{\piclt}[0]{\operatorname{\mathbf{Pic}}^{\rm loc-\tau}}
\newcommand{\piclp}[0]{\operatorname{\mathbf{Pic}}^{\rm loc-p}}
\newcommand{\pico}[0]{\operatorname{\mathbf{Pic}}^{\circ}}
\newcommand{\pict}[0]{\operatorname{\mathbf{Pic}}^{\rm \tau}}

\newcommand{\alb}[0]{\operatorname{Alb}}

\newcommand{\nsl}[0]{\operatorname{NS}^{\rm loc}}
\newcommand{\ns}[0]{\operatorname{NS}}



\newcommand{\gal}[0]{\operatorname{Gal}}
\newcommand{\cont}[0]{\operatorname{cont}}
\newcommand{\rank}[0]{\operatorname{rank}}
\newcommand{\mult}[0]{\operatorname{mult}}
\newcommand{\discrep}[0]{\operatorname{discrep}}
\newcommand{\totaldiscrep}[0]{\operatorname{totaldiscrep}}

\newcommand{\supp}[0]{\operatorname{Supp}}    
\newcommand{\red}[0]{\operatorname{red}}    
\newcommand{\codim}[0]{\operatorname{codim}}    
\newcommand{\im}[0]{\operatorname{im}}    
\newcommand{\proj}[0]{\operatorname{Proj}}    
\newcommand{\wt}[0]{\operatorname{wt}}    
\newcommand{\socle}[0]{\operatorname{socle}}    
\newcommand{\coker}[0]{\operatorname{coker}}    
\newcommand{\ext}[0]{\operatorname{Ext}}    
\newcommand{\Hom}[0]{\operatorname{Hom}}    
\newcommand{\trace}[0]{\operatorname{Trace}}  
\newcommand{\cent}[0]{\operatorname{center}}
\newcommand{\bs}[0]{\operatorname{Bs}}
\newcommand{\aut}[0]{\operatorname{Aut}}    
\newcommand{\specan}[0]{\operatorname{Specan}}    
\newcommand{\inter}[0]{\operatorname{Int}}    
\newcommand{\sing}[0]{\operatorname{Sing}}    
\newcommand{\ex}[0]{\operatorname{Ex}}    
\newcommand{\diff}[0]{\operatorname{Diff}}  
\newcommand{\diffg}[0]{\operatorname{Diff}^*}  
\newcommand{\tors}[0]{\operatorname{tors}}  
\newcommand{\tor}[0]{\operatorname{Tor}}  
\newcommand{\res}[0]{\operatorname{\mathcal R}} 
\newcommand{\disc}[0]{\operatorname{discr}} 
\newcommand{\jac}[0]{\operatorname{Jac}} 
\newcommand{\chow}[0]{\operatorname{Chow}}
\newcommand{\chr}[0]{\operatorname{char}}
\newcommand{\comp}[0]{\operatorname{Comp}}
\newcommand{\br}[0]{\operatorname{Br}}
\newcommand{\var}[0]{\operatorname{Var}}
\newcommand{\hilb}[0]{\operatorname{Hilb}}
\newcommand{\grass}[0]{\operatorname{Grass}}


\newcommand{\univ}[0]{\operatorname{Univ}}
\newcommand{\weil}[0]{\operatorname{Weil}}
\newcommand{\weilf}[0]{\operatorname{Weil_{tf}}}
\newcommand{\cl}[0]{\operatorname{Cl}}
\newcommand{\clns}[0]{\operatorname{\cl^{\rm ns}}}
\newcommand{\clo}[0]{\operatorname{\cl^{\circ}}}


\newcommand{\gr}[0]{\operatorname{gr}}
\newcommand{\len}[0]{\operatorname{length}}
\newcommand{\stab}[0]{\operatorname{stab}}

%\renewcommand{\conj}[0]{\operatorname{conj}} 
\newcommand{\ric}[0]{\operatorname{Ricci}}

\newcommand{\nec}[1]{\overline{NE}({#1})}

\newcommand{\rup}[1]{\lceil{#1}\rceil}
\newcommand{\rdown}[1]{\lfloor{#1}\rfloor}

\newcommand{\onto}[0]{\twoheadrightarrow}

\newcommand{\lcm}[0]{\operatorname{lcm}}
\newcommand{\wdeg}[0]{\operatorname{wdeg}}
\newcommand{\obs}[0]{\operatorname{obs}}
\newcommand{\ann}[0]{\operatorname{Ann}}
\newcommand{\diag}[0]{\operatorname{diag}}
\newcommand{\ord}[0]{\operatorname{ord}}
\newcommand{\mord}[0]{\operatorname{max-ord}}
\newcommand{\der}[0]{\operatorname{Der}}
\newcommand{\cosupp}[0]{\operatorname{cosupp}}
\newcommand{\simq}[0]{\sim_{\q}}
\newcommand{\simql}[0]{\sim_{\q,loc}}
\newcommand{\simr}[0]{\sim_{\r}}
\newcommand{\simrl}[0]{\sim_{\r,loc}}

\newcommand{\rest}[0]{\operatorname{rest}}
\newcommand{\isom}[0]{\operatorname{Isom}}


\newcommand{\uhom}[0]{\operatorname{\underline{Hom}}}  
\newcommand{\quot}[0]{\operatorname{Quot}}
\newcommand{\opic}[0]{\operatorname{\overline{Pic}}}  
\newcommand{\mor}[0]{\operatorname{Mor}} 
\newcommand{\depth}[0]{\operatorname{depth}} 
\newcommand{\tsum}[0]{\textstyle{\sum}} 
\newcommand{\tr}[0]{\operatorname{tr}} 

\newcommand{\bclq}[0]{\operatorname{B-Cl^{\q}}} 
\newcommand{\deffi}[0]{\operatorname{defficiency}} 
\newcommand{\bdiv}[0]{\operatorname{B-Div}} 
\newcommand{\wdiv}[0]{\operatorname{WDiv}} 
\newcommand{\cdiv}[0]{\operatorname{CDiv}} 
\newcommand{\pdiv}[0]{\operatorname{PDiv}} 


\newcommand{\sext}[0]{\operatorname{\mathcal{E}\!\it{xt}}}    
\newcommand{\shom}[0]{\operatorname{\mathcal{H}\!\it{om}}} 

\newcommand{\cond}[0]{\operatorname{cond}} 
\newcommand{\nonsnc}[0]{\operatorname{non-snc}} 
\newcommand{\snc}[0]{\operatorname{snc}} 
\newcommand{\homoth}[0]{\operatorname{Homothety}} 
\newcommand{\bir}[0]{\operatorname{Bir}} 
\newcommand{\nklt}[0]{\operatorname{nklt}} 
\newcommand{\ind}[0]{\operatorname{index}} 
\newcommand{\coeff}[0]{\operatorname{coeff}} 
\newcommand{\mld}[0]{\operatorname{mld}} 
\newcommand{\mldb}[0]{\overline{\operatorname{mld}}} 
\newcommand{\defor}[0]{\operatorname{Def}} 
\newcommand{\deforb}[0]{\operatorname{\mathbf{Def}}} 

\newcommand{\rc}[0]{\operatorname{RatCurve}} 

\newcommand{\mg}[0]{{\mathbb G}_m}

\newcommand{\nablax}[0]{\nabla_{\!\mathbf x}}


%\UseAMSsymbols

%\DeclareMathSymbol\subsetneq 2328
%\DeclareMathSymbol\onto 1310

%\DeclareMathSymbol\twoheadrightarrow 1310
%\DeclareMathSymbol\rtimes 226F
%\DeclareMathSymbol\ltimes 226E

\def\into{\DOTSB\lhook\joinrel\to}


\newcommand{\sF}{\mathcal{F}}
\newcommand{\sG}{\mathcal{G}}
\newcommand{\sM}{\mathcal{M}}
\newcommand{\sL}{\mathcal{L}}
\newcommand{\sN}{\mathcal{N}}
\newcommand{\sH}{\mathcal{H}}

\newcommand{\bolddot}{{\raisebox{.15em}{\ensuremath\centerdot}}}
\newcommand{\leteq}{\colon\!\!\!=}
\def\qis{\,{\simeq}_{qis}\,}

\def\loccoh#1.#2.#3.#4.{H^{#1}_{#2}(#3,#4)}

\DeclareMathAlphabet{\mathchanc}{OT1}{pzc}%
                                {m}{it}

\newcommand{\od}[1]{\ensuremath{\mathchanc{d}(\!{#1}\!)}}
%\newcommand{\sHom}[0]{{\mathchanc{Hom}}} 
%\newcommand{\sExt}[0]{{\mathchanc{Ext}}}  
\newcommand{\resto}{\big\vert_}
\newcommand{\rresto}{\right\vert_}


%\newcommand{\szabores}{Szab\'o-resolution }
\newcommand{\pairD}{D}

\newcommand{\gm}[0]{{\mathbb G}_m}
\newcommand{\src}[0]{\operatorname{Src}} 
\newcommand{\spr}[0]{\operatorname{Spr}} 
\newcommand{\simb}[0]{\stackrel{bir}{\sim}}
\newcommand{\sims}[0]{\stackrel{stab}{\sim}}

\newcommand{\hn}[0]{\operatorname{\mathbf{HN}}} 
\newcommand{\hb}[0]{\operatorname{\mathbf{HB}}} 

\newcommand{\age}[0]{\operatorname{age}}
\newcommand{\GL}{\mathrm{GL}}
\newcommand{\PGL}{\mathrm{PGL}}
\newcommand{\PSL}{\mathrm{PSL}}
\newcommand{\SU}{\mathrm{SU}}
\newcommand{\SL}{\mathrm{SL}}
\newcommand{\U}{\mathrm{U}}
\newcommand{\SO}{\mathrm{SO}}
\newcommand{\OO}{\mathrm{O}}
\newcommand{\PGO}{\mathrm{PGO}}



\newcommand{\OG}{\mathrm{OG}}
\newcommand{\OF}{\mathrm{OF}}
\newcommand{\G}{\mathrm{G}}


\newcommand{\PTC}{\mathrm{PTC}}

\newcommand{\sym}[0]{\operatorname{Sym}}


\newcommand{\shh}[0]{\operatorname{\mathcal{H}}} 


%\makeindex

%\input{z-preamb}


\usepackage[all]{xy}\xyoption{dvips}

\def\CC{\mathcal C}


\newcommand{\mdeg}[0]{\operatorname{Deg}}
\newcommand{\dres}[0]{\operatorname{\mathcal D\mathcal R}} 
\newcommand{\tprod}[0]{\textstyle{\prod}} 

\newcommand{\norm}[0]{\operatorname{norm}}
\newcommand{\ass}[0]{\operatorname{Ass}}
\newcommand{\link}[0]{\operatorname{link}}
\newcommand{\td}[0]{\operatorname{Td}} 
\newcommand{\ndiv}{\hspace{-4pt}\not|\hspace{2pt}} 

%\newcommand{\wcl}[0]{\operatorname{wcl}}
%\newcommand{\scl}[0]{\operatorname{scl}}



\newcommand{\scl}[2]{\operatorname{scl}\left(#1\leftarrow #2\right)}
\newcommand{\wcl}[2]{\operatorname{wcl}\left(#1\leftarrow #2\right)}

\newcommand{\emb}[0]{\operatorname{Emb}}
\newcommand{\reg}[0]{\operatorname{reg}}
\newcommand{\assc}[0]{\operatorname{\overline{ass}}}
\newcommand{\dx}[0]{\partial_x}
\newcommand{\dy}[0]{\partial_y}
\newcommand{\qg}[0]{\ensuremath{\mathrm{qG}}}
\newcommand{\embdim}[0]{\operatorname{embdim}}

\newcommand{\deform}[0]{\operatorname{Def}}


\newcommand{\maps}[0]{\operatorname{Map}}
\newcommand{\mapsd}[0]{\operatorname{Map}^{\rm deg}}
\newcommand{\mapsnd}[0]{\operatorname{Map}^{\rm nd}}

\newcommand{\mo}[0]{\operatorname{M}^{\rm \circ}}

\newcommand{\mdiv}[0]{\operatorname{MDiv}}
\newcommand{\mcl}[0]{\operatorname{MCl}}
\newcommand{\mdsh}[0]{\operatorname{MDSh}}
\newcommand{\dsh}[0]{\operatorname{DSh}}

\newcommand{\sima}[0]{\stackrel{\rm alg}{\sim}}


\newcommand{\lomega}{\tilde{\omega}}
\newcommand{\Lomega}{\tilde{\Omega}}




\newcommand{\tfs}{T\!f\!S_2}


\newcommand{\find}[0]{\operatorname{fx}}
\newcommand{\fdiv}[0]{\operatorname{div}}
\newcommand{\mfrd}[0]{\operatorname{\mathfrak{d}}}

\newcommand{\lcg}[0]{\operatorname{lcg}}

\newcommand{\lct}[0]{\operatorname{lct}}
\newcommand{\qgdef}[0]{\operatorname{Def_{qG}}}
\newcommand{\defsp}[0]{\operatorname{Def}}
%%xxx

\begin{document}
\bibliographystyle{amsalpha}

\title{KSB smoothings of surface pairs}
        \author{J\'anos Koll\'ar}

        \begin{abstract}  We describe KSB smoothings of log canonical surface  pairs  $(S, D)$, where $D$ is a reduced curve. In sharp contrast with the $D=\emptyset$ case, cyclic quotient  pairs  always have
          KSB smoothings, usually forming many irreducible components.
        \end{abstract}



        \maketitle

        \tableofcontents

At the boundary of the moduli space of smooth surfaces, we find surfaces with log canonical singularities. Quotient singularities form the largest---and usually most troublesome---subclass.

 The universal deformation space $\defsp(S)$ of a quotient singularity is quite complicated. The irreducible components were enumerated in \cite{ksb, MR1129040}; their number grows exponentially with the multiplicity. 
However, for the moduli of surfaces, we are interested only in those deformations of $S$ for which $K_S$ lifts to a $\q$-Cartier divisor; see Paragraph~\ref{ksb.defs.say.0}.
These are now called {\it KSB deformations,} replacing the name  {\it  \qg-deformations} used in \cite{ksb}. 

By \cite[3.10]{ksb},  quotient singularities either have no KSB smoothings, or 
the  KSB smoothings form a single,  easy to describe, irreducible component of the deformation space;  see Paragraph~\ref{ksb.defs.say}.




The aim of this note is to study the analogous question for  the moduli theory of  simple normal crossing   surface pairs. 
The singularities at the boundary are log canonical pairs $(S, D)$, where $S$ itself has quotient singularities whenever $D\neq \emptyset$. There are 3 possible types; see  Notation~\ref{conormal.notation} for their description using dual graphs.

We enumerate the irreducible components of $\defsp_{\rm KSB}(S,D)$---the {\it space of KSB deformations}---where we now require $K_S+D$ to lift to a $\q$-Cartier divisor; see Definitions~\ref{ksb.defs.pair.defn}--\ref{ksb.defs.pair.dd}.


\begin{thm} \label{cyc.nplt.enum.i} The irreducible components of
  $$
  \defsp_{\rm KSB}\bigl(\bullet \ - \ c_1  \ - \ c_2  \ - \ \cdots  \ - \  c_r  \ - \ \bullet\bigr)
  $$
  are in one-to-one correspondence with the 
irreducible components of
  $$
  \defsp\bigl(c_1  \ - \ c_2  \ - \ \cdots  \ - \  c_r\bigr),
  $$
  and, in each of them, a general deformation is  smooth.
\end{thm}

\begin{thm} \label{dih.nplt.enum.i}  The irreducible components of
  $$
  \defsp_{\rm KSB}\left(
  \begin{array}{ccc}
    & 2 &\\
     & | &\\
    2  \ - \!\!\!\!   &c_1& \!\!\!\!   - \ c_2  \ - \ \cdots  \ - \  c_r  \ - \ \bullet
  \end{array}
  \right)
  $$
  are in one-to-one correspondence with the 
irreducible components of
  $$
\defsp\bigl(c_2  \ - \ \cdots  \ - \  c_r\bigr),
\eqno{\mbox{(note that  $c_1$ is omitted)}}
  $$
   and, in each of them, a general deformation has 2 singularities,  which are   $(2 \ - \ \bullet)$. 
\end{thm}

\begin{prop} \label{cyc.plt.rig.i}
  The singularities $( c_1  \ - \ c_2  \ - \ \cdots  \ - \  c_r\ - \ \bullet)$ are KSB rigid.
  \end{prop}



  
\begin{rems}{\ }
\begin{enumerate}
\item  The dimensions of the irreducible components of $\defsp_{\rm KSB}(S, D)$ are computed in (\ref{main.cyclic.thm}) for the cyclic cases, and in (\ref{main.cyclic.thm.d}) for the dihedral  cases.
\item  For the cyclic quotients  in Theorem~\ref{cyc.nplt.enum.i},
  the image of
  $$
  \defsp_{\rm KSB}\bigl(\bullet \ - \ c_1 \ - \ \cdots  \ - \  c_r  \ - \ \bullet\bigr)\ \to\ \defsp\bigl(c_1  \ - \cdots  \ - \  c_r\bigr)
  $$
  is nowhere dense, with a few exceptions; see (\ref{dim.drop.cyc.say}). 
  \item For quotient singularities $S$, \cite[3.14]{ksb} gives an algorithm to
enumerate the irreducible components of $\defsp(S)$, though in practice this can be very cumbersome if there are many exceptional curves. For cyclic quotients $S$ as in Theorem~\ref{cyc.nplt.enum.i}, \cite{MR1129040} gives a  better method.  As a consequence, \cite{MR1129040} shows that if $S$ has  multiplicity $m=m(S)$, then $\defsp(S)$ has at most  $\tfrac1{m-1}\tbinom{2(m-2)}{m-2}$
  irreducible components, and equality holds for `most' cyclic quotients of a given multiplicity. Note that $m(S)=2+\sum (c_i-2)$. 
    \end{enumerate}
\end{rems}

{\it Other results.} 
If we  look at pairs $(S, dD)$ for  $d<1$, the situation changes dramatically. For the pairs $(S, dD)$ as in Theorems~\ref{cyc.nplt.enum.i}--\ref{dih.nplt.enum.i}, there are either no  KSB deformations for any $d$, or a single
irreducible component for every $d$.
For almost all of the pairs $(S, dD)$ as in Proposition~\ref{cyc.plt.rig.i}, there are 
 either no  KSB deformations for any $d$, or a single
 irreducible component for a unique value of  $d$, see (\ref{4.0.smooth.exmp}). For $d\in (\frac56,1]$ we give a complete answer   in Section~\ref{0-1.sec}.
 For $d\in [\frac12,1]$, all lc pairs $(S, dD)$ are listed   in \cite[Sec.3.3]{kk-singbook}. Using it, 
  the method should give a full answer for  $d\in [\frac12,1]$, but I checked  only  some examples. Exceptional examples with $d=\frac12$ and  KSB smoothings are given in (\ref{1/2.2.exmp}).


\medskip




{\it Sketch of the proofs.}
Let $(S, D)$ be a pair as in
Theorems~\ref{cyc.nplt.enum.i}--\ref{dih.nplt.enum.i}, fix  an irreducible component of
  $\defsp_{\rm KSB}(S, D)$,  and let
$g:({\mathbf S}, {\mathbf D})\to  (0,\dd)$
be a general 1-parameter deformation in it. That is,  the fiber over $0$ is
 $({\mathbf S}, {\mathbf D})_0\cong (S, D)$, and $\dd$ denotes either a complex disc or the germ of a smooth curve.

 We prove in (\ref{smoothing.prop}) that the general fiber of $g$ has at worst $A_1$ singularities.  Thus, as we recall in (\ref{cmod.ksb.thm}), by \cite[3.5]{ksb} there is a 
proper, birational morphisms $\pi:{\mathbf S}_P\to {\mathbf S}$ such that
 \begin{itemize}
 \item $K_{{\mathbf S}_P}$ is $\pi$-ample, and
   \item the central fiber 
 $S_P$ has only  Du~Val and $\a^2/\tfrac1{rn^2}(1, arn{-}1)$  singularities.
 \end{itemize}
 (The latter are called T-singularities  (\ref{ksb.defs.say}.1), and
 $S_P\to S$ is a 
   P-modification, see   Definition~\ref{P+M.defn}).
 
 Let ${\mathbf D}_P$ denote the birational transform of ${\mathbf D}$ on
 ${\mathbf S}_P$.  Then
 $K_{{\mathbf S}_P}+{\mathbf D}_P\simq \pi^*\bigl(K_{{\mathbf S}}+{\mathbf D}\bigr)$, hence 
 $g\circ\pi:({\mathbf S}_P+{\mathbf D}_P)\to \dd$
 is a KSB deformation of $S_P$ and also of   $(S_p, {\mathbf D}_P|_{S_P})$;
 we call these doubly KSB deformations  in Definition~\ref{2-KSB.defn}.
 We show in (\ref{K+D.KSB.P.thm}) that ${\mathbf D}_P|_{S_P}=D_P+E_P$, where
 $E_P\subset S_P$ is the reduced exceptional divisor.

 A well known argument (\ref{main.MP.old.thm}) then  gives  a natural morphism
$\defsp_{\rm KSB}(S_P, D_P+E_P) \to \defsp_{\rm KSB}(S,D)$, which is 
      finite and  birational  onto our irreducible component.
   
In the cyclic case of Theorem~\ref{cyc.nplt.enum.i}, a quick argument shows that every P-modification leads to an  irreducible component of 
$\defsp_{\rm KSB}(S,D)$, see (\ref{main.cyclic.thm}).

In the dihedral  case of Theorem~\ref{dih.nplt.enum.i},  the irreducible components of $\defsp(S)$ are enumerate in (\ref{dih.P.mod.all}).
Then in (\ref{DP.char.lem}--\ref{dih.P.mod.all.rem})  we describe which P-modification lead to an  irreducible component of 
$\defsp_{\rm KSB}(S,D)$.    Putting these together gives 
Theorem~\ref{dih.nplt.enum.i}.

Proposition~\ref{cyc.plt.rig.i}  directly follows from \cite[2.23]{k-modbook}, see (\ref{ord.of.K+C.say}.1). \qed

 


\medskip

{\it Infinitesimal computations.}  Usually the space of KSB deformations $\defsp_{\rm KSB}(S, D)$ is not reduced. For cyclic pairs, the method of  \cite{k-alt} is used to  determine its tangent space in Section~\ref{first.ord.sec}, but  a complete description is known only in a  few cases. 



  \medskip

{\it Other applications.}
Many of the results  apply to more general rational singularities  $S$.  However,  not   every irreducible component of $\defsp(S)$
 is obtained from a P-modification.


The conjectures in \cite[Sec.6]{k-etc} ask whether every irreducible component of $\defsp(S)$ is obtained in a similar way, using a  notion of P-modification that is more general than the one in Definition~\ref{P+M.defn}.  A positive answer  has been known for quotient singularities  \cite[3.9]{ksb} and  for  points of multiplicity $\leq 4$ \cite{MR1135880, MR1094706}.  The recent papers \cite{park2022deformations, jeon2023deformations} develop a method to prove the conjecture in many new cases. In \cite{park2022deformations} this is illustrated by the
$W(p,q,r)$ series of singularities, but their method  applies more broadly.  (The $W(p,q,r)$ series was  discovered by Wahl around 1980,  see \cite{MR4270358} for a recent survey.)  In \cite{jeon2023deformations} the conjecture is proved for most weighted homogeneous singularities.


Especially for the M-modification version  as in \cite{be-ch} or (\ref{main.MP.old.thm}), the
singularities with a rational homology disc smoothing---classified by \cite{MR2843099}---may form the natural class to work with.

\medskip

{\it Higher dimensions.} For an lc pair $(X, D)$ of arbitrary dimension with reduced $D$, \cite{k-db3} shows that a flat deformation is a KSB  deformation iff the KSB condition (\ref{ksb.defs.pair.defn})  is satisfied by a  general surface section.
  Thus our results give  information not only for surface pairs, but for higher dimensions as well.



\begin{ack}  I thank  K.~Altmann, P.~Hacking, H.~Hauser and J.~Stevens  for many   useful comments and references.
Partial  financial support    was provided  by  the NSF under grant number
DMS-1901855.
\end{ack}


 

\section{Quotient singularities}


The classification of surface quotient singularities and  their dual graphs are given in \cite{briesk}.
We need detailed information about the cyclic (type $A$) and dihedral (type $D$) cases; we recall these and fix our notation.
The tetrahedral, octahedral and icosahedral quotients (type $E$) will appear only in examples.

\begin{notation}
We work over the complex numbers and let $\a^2$ denote the germ of ${\mathbb A}^2$ at the origin. For our purposes, we may work with a complex analytic germ,
the spectrum of $\c[[x,y]]$, or the  spectrum of the Henselisation of
$\c[x,y]_{(x,y)}$.
For most situations in this paper one can choose  global coordinates, so we can even work with   $0\in {\mathbb A}^2$.

$S_{n,q}$ and $S^d_{n,q}$ will denote the cyclic and dihedral
quotients as in Notation~\ref{cyclic.basic.not} and \ref{dih.basic.not}.
The curves  $B_{n,q}\subset S_{n,q}$, $D_{n,q}\subset S_{n,q}$ and $D^d_{n,q}\subset S^d_{n,q}$ are defined in (\ref{conormal.notation}).

The minimal resolution of a surface $S$ is denoted by $\mu:S^{\rm m}\to S$.
For a curve $D\subset S$, its birational transform on $S^{\rm m}$ is denoted by $D^{\rm m}$.
\end{notation}

\begin{notation}[Cyclic quotients]\label{cyclic.basic.not}
  We write $S_{n,q}:=\a^2/\frac1{n}(1,q)$
  where the group action is $(x, y)\mapsto (\epsilon x, \epsilon^q y)$
  for some primitive $n$th root of unity $\epsilon$. The action is free outside the origin iff $(n, q)=1$.   In general, if $n=n_1(n,q), q=q_1(n,q)$, then
  $$
  \a^2/\tfrac1{n}(1,q)\cong \a^2/\tfrac1{n_1}(1,q_1).
  $$
  Let $q'$ denote the multiplicative inverse of $q$  modulo $n$. Then
  $S_{n,q}\cong  S_{n,q'}$;  the isomorphism interchanges the coordinates.

  For the dual graph of the minimal resolution, we use $c_i:=-(C_i^2)$, the negative of the self-intersection, to denote the vertex corresponding to the curve $C_i\subset S_{n,q}^{\rm m}$.
   For $S_{n,q}$ the dual graph is
  $$
 {\mathcal D}(S_{n,q}) =\quad  c_1\ - \ c_2 \ - \cdots \  - c_r,
  $$
    where the $c_i$ are obtained from the continued fraction expansion
    $$
    \frac{n}{q}=c_1-\cfrac{1}{c_2-\cfrac{1}{c_3-\cdots}}
    $$
    We use $[c_1, \dots, c_r]$ to denote this continued fraction, and
    $S(c_1, \dots, c_r)\cong S_{n,q}$ the resulting singularity.
      \end{notation}

  
\begin{say}[Resolution of cyclic quotients]  \label{res.of.cyclic.say}
An explicit construction of the minimal resolution is in \cite{jung};
  see \cite{reid-cyclic} for a very accessible treatment.
  Here we use an inductive  method as in \cite[2.31]{k-res}.

  For ${\mathbb A}^2_{xy}/\frac1{n}(1,q)$, the quotient of the blow-up
  of the ideal sheaf $(x^q, y)$ gives a proper, birational morphism
  $\pi_1:S_1\to  {\mathbb A}^2_{xy}/\frac1{n}(1,q)$. It is covered by 2 charts
  \begin{enumerate}
  \item (Singular chart)  ${\mathbb A}^2_{x_1y_1}/\frac1{q}(1,-n)$, where
    $x=x_1y_1^{1/n}, y=y_1^{q/n}$.
    \item (Smooth chart)  ${\mathbb A}^2_{x_2y_2}$, where
      $x=x_2^{1/n}, y=y_2x_2^{q/n}$.
  \end{enumerate}
   Iterating this blow-up gives the minimal resolution.
   The $\pi_1$-exceptional curve $E_1$ is $(y_1=0)$ (resp.\ $(x_2=0)$).
   Thus the extended dual graph is
   $$
   (x-\mbox{axis})\ -\   c_1\ - \ c_2 \ - \cdots \  - c_r\ -\ (y-\mbox{axis}).
   \eqno{(\ref{res.of.cyclic.say}.3)}
  $$
\end{say}


\begin{notation}[Dihedral quotients]\label{dih.basic.not}
  Let $S^d_{n,q}$ denote the singularity whose
  minimal resolution dual graph  is
  $$
  \begin{array}{ccc}
    & 2 &\\
     & | &\\
    2  \ - \!\!\!\!   &c_1& \!\!\!\!   - \ c_2  \ - \ \cdots  \ - \  c_r 
  \end{array}
  $$
%%     $$
%% \xymatrix{%
%%    &2&  \\
%%   2  \ar@{-}[r] & c_1  \ar@{-}[u] \ar@{-}[r] & \cdots \ar@{-}[r] & c_r  }
%% \eqno{(\ref{dih.basic.not}.1)}
%% $$
where $\frac{n}{q}=[c_1, \dots, c_r]$.  The curves marked $2$ are denoted by $C'_0, C''_0$. 
We also use the notation
$$
S(2^2, c_1,\dots, c_r):=S^d_{n,q}.
$$
One usually assumes $r\geq 2$, though later it will be convenient to allow $r=1$.

If we contract the curves  $C'_0, C''_0$ and $C_2, \dots, C_r$, we get
$\bar S^d_{n,q}\to S^d_{n,q}$. There is 
a single exceptional curve $\cong \p^1$ with two $A_1$  points on  it, plus the singularity coming from
 $S(c_2, \dots, c_r)$. Thus $\bar S^d_{n,q}$ has a double cover, ramified only at the  $A_1$  points. The corresponding  double cover of $S^d_{n,q}$ is a cyclic   quotient singularity, 
whose dual graph is
$$
c_r \ - \ \cdots\  - \ c_2\  - \ \tilde c_1\  - \ c_2   \ - \ \cdots\  - \ c_r
\eqno{(\ref{dih.basic.not}.2)}
$$
 where $ \tilde c_1= 2(c_1-1)$.
   %% $$
   %%  \xymatrix{%
   %%    \bullet \ar@{-}[r] & c_1 \ \cdots \  c_{r-1} \ar@{-}[r] & \tilde c_r\ar@{-}[r]   & c_{r-1}\ \cdots\ c_1 \ar@{-}[r]   & \bullet }
   %%  $$
We claim that   this is the quotient singularity
$$
S_{N,Q}:=\a^2/\tfrac1{2q(n-q)}\bigl(1, 2n'(n-q)+1\bigr), \qtq{where} qq'=nn'+1,
\eqno{(\ref{dih.basic.not}.3)}
$$
giving the correspondence
$$
\begin{array}{cll}
  (n,q) & \mapsto & \bigl(N=2q(n-q),\ Q=2n'(n-q)+1\bigr), \qtq{and}\\
  (N,Q) & \mapsto & \bigl(n= q+ \tfrac12(N,Q-1),\   q=N/(N,Q-1)\bigr).
\end{array}
\eqno{(\ref{dih.basic.not}.4)}
$$
(So  $N$ is even and $Q^2\equiv 1\mod N$.)
To see this, start with the singularity (\ref{dih.basic.not}.3) and blow up the origin.
In the $(x, \frac{y}{x})$-chart  we get
$$
{\mathbb A}^2/\tfrac1{2q(n-q)}\bigl(1, 2n'(n-q)\bigr)\cong 
{\mathbb A}^2/\tfrac1{q}\bigl(1, n'\bigr).
$$
Next, write
$\tfrac{q}{m}:=[c_2,\dots, c_r],$ so  $n=c_1-\tfrac{m}{q}$,
and note that   
$nn'\equiv -1 \mod q$, so $mn' \equiv 1 \mod q$.
Therefore
$$
\a^2/\tfrac1{q}\bigl(1, n'\bigr)\cong
\a^2/\tfrac1{q}\bigl(1, m\bigr)
$$
which gives $S(c_2,\dots, c_r)$.

The exceptional curve of the blow-up is   pointwise fixed by a subgroup of order $2(n-q)$. Taking the corresponding quotient we have  a smooth surface and the self-intersection of the curve is $-2(n-q)$. Then we have a $\frac1{q}$-action, which changes the  self-intersection to $-2(n-q)/q=-2(n/q)+2$. Then we resolve the  $\a^2/\tfrac1{q}\bigl(1, m\bigr)$ singularities. The self-intersection becomes   $-2\rup{n/q}+2=-2c_1+2$, as claimed.



\end{notation}
%%   $$
%%       \xymatrix{%
%%       \bar S^{\kappa}  \ar[r]^{\mu^{\kappa}} \ar[d]_{\pi^m}   & S^{\kappa}  \ar[d]^{\pi^{\kappa}} \\
%% S^{\rm m}\ar@{.>}[ur]^{\tau}  \ar[r]^{\mu}   &  S
%%     }
%%     $$



\begin{say}[KSB deformations]\label{ksb.defs.say.0}
  Let $X_0$ be a surface with quotient (more generally log canonical) singularities only. A {\it KSB deformation} of $X_0$ 
  is a flat morphism to a local scheme $g: X\to (0,B)$ such that
  $X_0\cong g^{-1}(0)$, and
  \begin{enumerate}
  \item $\omega_{X/B}^{[m]}$ is flat over $B$, and  commutes with all base changes $(0,B')\to (0, B)$ for every $m$.
  \end{enumerate}
  If $B$ is reduced, then, by \cite[2.3 and 4.7]{k-modbook} this is equivalent to 
\begin{enumerate}\setcounter{enumi}{1}
  \item $\omega_{X/B}^{[m]}$ is locally free for some  $m>0$.
\end{enumerate}
The notion was introduced in \cite{ksb}, and called {\it  $\qg$-deformation} there. Note also that the 2 versions are not equivalent already over
$B=\spec \c[\epsilon]$; see \cite{k-alt}. 
\end{say}

\begin{say}[KSB smoothings]\label{ksb.defs.say}
  By \cite[3.10]{ksb}, quotient singularities that admit a KSB smoothing---frequently called {\it T-singularities---}are either Du~Val or of the form
  $$
  \a^2_{uv}/\tfrac1{rn^2}(1, arn-1),\qtq{where $(a, n)=1$ and $n>1$.}
  \eqno{(\ref{ksb.defs.say}.1)}
  $$
  More generally, these are the singularities that have a KSB deformation with Du~Val generic fiber.
  (There are a few more cases that have  nontrivial KSB deformations, but no KSB smoothings; see \cite[2.29]{k-modbook}.)
  By \cite{MR83h:14029}, the dual resolution graphs are obtained as follows.
  We start with
  $$
   4 
   \qtq{or}
         3 \ - \ 2 \ - \ \cdots \ - \ 2\ - \ 3,
    $$
    and successively apply the operation
    $$
         c_1 \ - \ \cdots \ - \   c_{r-1}  \ - \   c_r 
  \quad  \mapsto \quad
      2  \ - \   c_1  \ - \ \cdots \ - \   c_{r-1}  \ - \   (c_r+1), 
    $$
      or its symmetric version.

      Starting with $(4)$ gives the $\a^2_{uv}/\tfrac1{n^2}(1, arn-1)$
      singularities, and starting with $(3 \ - \ 2 \ - \ \cdots \ - \ 2\ - \ 3)$
      (with $r{-}2$ curves marked $2$) gives the
      $\a^2_{uv}/\tfrac1{rn^2}(1, arn-1)$ singularities.

  
    Taking the quotient of (\ref{ksb.defs.say}.1)
    by the subgroup of order $rn$ shows that
   $$
    \a^2_{uv}/\tfrac1{rn^2}(1, arn-1)\cong (xy=z^{rn})/\tfrac1{n}(1, -1, a),
    \eqno{(\ref{ksb.defs.say}.2)}
  $$
  where $x=u^{rn}, y=v^{rn}, z=uv$. 
  Using the second representation, the universal KSB deformation is given by
  $$
  \bigl(xy=z^{rn}+\tsum_{i=0}^{r-1}t_iz^{in}\bigr)/\tfrac1{n}(1, -1, a, {\mathbf 0}).
   \eqno{(\ref{ksb.defs.say}.3)}
   $$
   Thus the KSB deformation space is smooth and has dimension $r$.
   

   For $r=1$ we get a 1-dimensional KSB deformation space. 
   As observed by \cite{be-ch}, the $r>1$ cases can be reduced to the $r=1$ cases as follows.
Change the deformation (\ref{ksb.defs.say}.3) to the factored form
  $$
  \bigl(xy=\tprod_{i=1}^{r}(z^{n}-s_i)\bigr)/\tfrac1{n}(1, -1, a, {\mathbf 0}),
  %% \eqno{(\ref{ksb.defs.say}.4)}
  $$
and  repeatedly blow up the divisors $(x=z^{n}-s_i=0)$
  to obtain a small, cerpant modification, which is a flat KSB deformation of its central fiber. The central fiber is a modification of 
  $\a^2/\tfrac1{rn^2}(1, arn{-}1)$  that has $r$ singularities of the form
$\a^2/\tfrac1{n^2}(1, an{-}1)$, connected by $r-1$ curves whose birational transforms are $(-1)$-curves on the minimal resolution of the central fiber. Using Notation~\ref{boxed.nota}, the dual graph  is the following.
$$
  {\boxed{\tfrac1{n^2}(1, an{-}1)}-}\ 1\ {-\boxed{\tfrac1{n^2}(1, an{-}1)}-}\  \cdots \ {-\boxed{\tfrac1{n^2}(1, an{-}1)}-} \ 1\ {-\boxed{\tfrac1{n^2}(1, an{-}1)}}
   \eqno{(\ref{ksb.defs.say}.4)}
   $$
   This construction leads to the notion of M-modifications in (\ref{P+M.defn}).

   Note that, as we repeatedly contract $(-1)$-curves in  (\ref{ksb.defs.say}.4), we never contract the curves on the two ends. These show the following.
   \medskip

   {\it Claim \ref{ksb.defs.say}.5.}  The minimal resolution of
   $\a^2/\tfrac1{rn^2}(1, arn{-}1)$ and the minimal resolution of (\ref{ksb.defs.say}.4) are isomorphic along the birational transforms of the $x$ and $y$ axes. \qed
   \medskip




   {\it Claim \ref{ksb.defs.say}.6.}
 The map  $\a^n_{\mathbf s}\to \a^n_{\mathbf t}$, given by the elementary symmetric functions, is a Galois cover of degree $n!$. \qed
  \end{say}



  \begin{notation} \label{boxed.nota} For $\frac{n}{q}=[c_1,\dots, c_r]$ we use
    $$
    a_1\ {-\boxed{\tfrac{n}{q}}-}\ a_2   \qtq{resp.}
    a_1\ {-\boxed{\tfrac{n}{q}}}
    $$
    to denote 2 (resp.\ 1) curves through a singular point of type $\a^2/\frac1{n}(1,q)$,  whose resolution dual graph is
    $$
    a_1\ -\ c_1 \ -\  \ \cdots \ - \ c_r \ -\ a_2
    \qtq{resp.}
     a_1\ -\ c_1 \ -\  \ \cdots \ - \ c_r.
    $$
    \end{notation}
    

    
\section{Surface pairs with reduced boundary}

\begin{notation} \label{conormal.notation}
  Let  $(S, D)$ be an lc surface pair, where $D$ is a reduced curve.
 There are 3 possible local normal forms at points  $p\in D\subset S$; see for example \cite[Sec.3.3]{kk-singbook}. 
\begin{enumerate}
\item   {\it (Cyclic, plt)}
  $(S_{n,q}, B_{n,q}):=\bigl(\a^2_{xy}, (x=0)\bigr)/\tfrac1{n}(1,q)$, where $(n,q)=1$.  The minimal resolution dual graph is the following, where $\frac{n}{q}=[c_1,\dots, c_r]$ and $\bullet$ denotes the birational transform of $B_{n,q}$. 
  $$
  c_1  \ - \ c_2  \ - \ \cdots  \ - \  c_r \ - \ \bullet
     $$
\item
   {\it (Cyclic, non-plt)}
   $(S_{n,q}, D_{n,q}):=\bigl(\a^2_{xy}, (xy=0)\bigr)/\tfrac1{n}(1,q)$, where $(n,q)=1$.
     The minimal resolution dual graph is
   $$
  \bullet \ - \  c_1  \ - \ c_2  \ - \ \cdots  \ - \  c_r \ - \ \bullet
  $$
  Note that $D_{n,n{-}1}$ is the Du~Val singularity  $(uv=w^n)$. These are the only ones for which $D_{n,q}$ is Cartier and they behave exceptionally in many respects.
\item
  {\it (Dihedral)}  The pair  $(S^d_{n,q}, D^d_{n,q})$ has dual graph.
  $$
  \begin{array}{ccc}
    & 2 &\\
     & | &\\
    2  \ - \!\!\!\!   &c_1& \!\!\!\!   - \ c_2  \ - \ \cdots  \ - \  c_r  \ - \ \bullet
  \end{array}
%% \xymatrix{%
%%    &2&&  \\
%%     2  \ar@{-}[r] & c_1  \ar@{-}[u] \ar@{-}[r] & \cdots \ar@{-}[r] & c_r  \ar@{-}[r] &  \bullet}
$$
The curve $D^d_{n,q}$ is irreducible and smooth.
As in (\ref{dih.basic.not}.3), $S^d_{n,q}$  is also obtained as the quotient of
$$
(S_{N,Q}, D_{N,Q}):=\bigl(\a^2, (xy=0)\bigr)/\tfrac1{2q(n-q)}\bigl(1, 2n'(n-q)+1\bigr)
$$
by an involution.
For $r=1$   the singularity $S^d_{n,q}$ is a cyclic quotient, but the curve $D^d_{n,q}$ is different from the cyclic cases.
 \end{enumerate}
\end{notation}




\begin{defn}[KSB deformations of pairs]\label{ksb.defs.pair.defn}
  Let $(X_0, D_0)$ be a normal, log canonical surface pair. A {\it KSB deformation} of $(X_0, D_0)$ 
  is a flat morphism to a local scheme $g: X\to (0,B)$, plus a divisor $D\subset X$  that is Cartier at the smooth points of $X_0$,   such that
  \begin{enumerate}
  \item $X_0\cong g^{-1}(0)$,
    \item $D$ is flat over $B$  and $D_0:=D|_{X_0}$,
  \item $\omega_{X/B}^{[m]}(mD)$ is flat over $B$, and  commutes with all base changes $(0,B')\to (0, B)$ for every $m$.
  \end{enumerate}
  If $B$ is reduced, then, by \cite[2.3 and 4.7]{k-modbook} the latter  is equivalent to 
\begin{enumerate}\setcounter{enumi}{3}
  \item $\omega_{X/B}^{[m]}(mD)$ is locally free for some  $m>0$.
\end{enumerate}
\end{defn}


\begin{say}[Deformation spaces]\label{ksb.defs.pair.dd}
  Let $Y$ be an affine variety (or Stein space) with isolated singularities and
  $p:X\to Y$  a proper, birational morphism such that $\ex(p)$ is proper.
  Then a universal deformation space $\defsp(X)$ exists (and is finite dimensional); see \cite{MR54:10261, bing}.
  We get $\defsp(Y)$ using the identity map $Y=Y$.
  
  If $Z\subset X$ is a closed subvariety with isolated singularities, then $\defsp(X, Z)$ exists.

  If $x\in X$ is a point, then $\defsp(x, X, Z)$ denotes the deformation space  of a small enough  neighborhood of $x\in U\subset X$.   

  If $D\subset X$ is a divisor  with isolated singularities, 
all KSB deformations give a  functor, that has a universal deformation space, denoted by
$$
\defsp_{\rm KSB}(X,D).
\eqno{(\ref{ksb.defs.pair.dd}.1)}
$$
As the simplest example, adding a divisor $D$ to (\ref{ksb.defs.say}.2)
gives a  smooth deformation component of
$\defsp_{\rm KSB}\bigl(\a^2_{uv}, (uv=0)\bigr)/\frac1{rn^2}(1, ran-1)$
$$
  \bigl((xy=z^{rn}+\tsum_{i=0}^{r-1}t_iz^{in}), (z=0) \bigr)/\tfrac1{n}(1, -1, a, {\mathbf 0}).
   \eqno{(\ref{ksb.defs.pair.dd}.2)}
   $$
Our main aim is to describe the reduced structure of $\defsp_{\rm KSB}(X,D)$ for log canonical surface pairs.
\end{say}

As in  \cite{ksb} or \cite[2.23]{k-modbook}, there is a simplifying step if
$K_{X}+D$ is not Cartier. We study this possibility next.



\begin{say}[The order of $K_S+D$ in the class group]\label{ord.of.K+C.say}
  We consider separately the  3
  cases in  (\ref{conormal.notation}.1--3).
\begin{enumerate}
\item  For (\ref{conormal.notation}.1), the group acts faithully on 
$\frac{dx\wedge dy}{x}$. Thus  $(\a^2_{xy}, (x=0)\bigr)$ is the index 1 cover of $(S_{n,q}, B_{n,q})$; see \cite[2.49]{kk-singbook}. Therefore, as in \cite[2.23]{k-modbook},
every KSB deformation of $(S_{n,q}, B_{n,q})$ is induced by a KSB deformation of  $(\a^2_{xy}, (x=0)\bigr)$. The latter is rigid, so every
KSB deformation of $(S_{n,q}, B_{n,q})$ is trivial.
\item  For (\ref{conormal.notation}.2),
  the rational 2 form
  $\frac{dx\wedge dy}{xy}$ is invariant under the group action, hence descends to a generating section of  $\omega_{S_{n,q}}(D_{n,q})$.  Thus
  $\omega_{S_{n,q}}(D_{n,q})$ is locally free and $K_S+D$ is Cartier.
  Thus the method of \cite[2.23]{k-modbook} does not give any information about the  KSB deformations.

  
  
  
\item  For (\ref{conormal.notation}.3),  the rational 2 form
  $\frac{dx\wedge dy}{xy}$ is invariant under the cyclic subgroup of order $2q(n-q)$, but anti-invariant under the remaining involution. Thus only its tensor square descends, so
  $K_S+D$ is 2-torsion in the local Picard group.
  The corresponding duble cover is $(S_{N,Q}, D_{N,Q})$ as in (\ref{conormal.notation}.3).

  As in \cite[2.23]{k-modbook},
  every KSB deformation of $(S^d_{n,q}, D^d_{n,q})$ is induced by a KSB deformation of  $(S_{N,Q}, D_{N,Q})$. Thus, in principle, this reduces us to the previous case. We give a more direct description of all KSB deformations in (\ref{dih.P.mod.all}).
 
\end{enumerate}

Pulling $K_S+D$ back to the minimal resolution shows that
\begin{enumerate}\setcounter{enumi}{3}
 \item  $K_{S^{\rm m}}+D^{\rm m}+\sum_{i=1}^r C_i\sim 0$ in case (\ref{conormal.notation}.2), and
\item   $2\bigl(K_{S^{\rm m}}+D^{\rm m}+\sum_{i=1}^r C_i\bigr)+(C'_0+C''_0) \sim 0$  case (\ref{conormal.notation}.3).  
\end{enumerate}
\end{say}


\section{Q-modifications}


P-modifications were introduced in \cite{ksb} (under the name P-resolution) to describe   smoothings of quotient singularities. A variant, using M-modifications is developed in  \cite{be-ch}.  To handle arbitrary deformations of pairs we introduce 
Q-modifications.  These are actually very natural objects from the point of view of birational geometry, and many basic results on P- and M-modifications
hold for Q-modifications. 

\begin{defn}\label{Q-mod.defn}
  Let $(s,S)$ be a rational surface singularity. A proper, birational morphism $\pi:S_Q\to S$ is a {\it Q-modification}
if
\begin{enumerate}
\item $K_{S_Q}$ is $\pi$-nef, and
\item $S_Q$ has only quotient singularities. (So `Q' for quotient.)
\end{enumerate}
\end{defn}


\begin{notation}\label{P.res.ample.say}
Let $(0,S)$ be a rational singularity with minimal resolution
$\mu:S^{\rm m}\to S$. Let $\pi:\bar S\to S$ be a proper, birational morphism.
Assume that  $\bar S$ is normal with  minimal resolution
 $\bar\mu:\bar S^{\rm m} \to \bar S$.
Then  $\bar S^{\rm m}$ dominates $S^{\rm m}$, giving a commutative diagram
  $$
    \xymatrix{%
      \bar S^{\rm m}  \ar[r]^{\bar\mu} \ar[d]_{\pi^m}   & \bar S  \ar[d]^{\pi} \\
S^{\rm m}\ar@{.>}[ur]^{\tau}  \ar[r]^{\mu}   &  S
    }
    \eqno{(\ref{P.res.ample.say}.1)}
 $$   
    We can write
    $$
    K_{\bar S^{\rm m}}\simq  \bar\mu^* K_{\bar S}-\bar\Delta,\qtq{and}
    K_{\bar S^{\rm m}}\sim (\pi^m)^*K_{S^{\rm m}}+F,
    \eqno{(\ref{P.res.ample.say}.2)}
    $$
    where $\bar\Delta$ is effective with $\supp \bar\Delta\subset \ex(\bar\mu)$,
    and $F$ is effective  with $\supp F = \ex(\pi^m)$.

        Let $\bar A^{\rm m}\subset \bar S^{\rm m}$ be an irreducible  curve that is exceptional over $S$, with  images  $A^{\rm m}\subset  S^{\rm m}$ and $\bar A\subset \bar S$. (We use $0$ if the image is  a point.)
    Then (\ref{P.res.ample.say}.2) gives that
$$
      (\bar A\cdot K_{\bar S})=
      (A^{\rm m}\cdot K_{S^{\rm m}})+\bigl(\bar A^{\rm m}\cdot F\bigr)+
      \bigl(\bar A^{\rm m}\cdot \bar \Delta\bigr).
      \eqno{(\ref{P.res.ample.say}.3)}
      $$
      From this we conclude the following.
      \begin{enumerate}\setcounter{enumi}{3}
      \item  If $\bar A\neq\emptyset$ and $ A^{\rm m}\neq\emptyset$, then
        $(\bar A\cdot K_{\bar S})\geq 0$, with equalty holding iff
        $(A^{\rm m}\cdot K_{S^{\rm m}})=0$ and $\tau$ is an isomorphism along
        $A^{\rm m}$.
      \item If $ A^{\rm m}=\emptyset$, then
        $ (\bar A\cdot K_{\bar S})$ depends only on the intersection numbers of the curves in $\ex(\bar\mu)\cup \ex(\pi^m)$.
      \end{enumerate}
\end{notation}

\begin{lem}\label{antieff.lem} Using the notation of  (\ref{P.res.ample.say}), assume that $K_{\bar S}$  is $\pi$-nef. 
  Then 
    $-K_{\bar S^{\rm m}}$ is $\q$-linearly equivalent to an effective linear combination of curves, that are exceptional for $\bar S^{\rm m}\to S$.
\end{lem}

Proof. By (\ref{P.res.ample.say}.2) $ K_{\bar S^{\rm m}}\simq  \bar\mu^* K_{\bar S}-\bar\Delta$ where $\Delta $ is effective. Since $K_{\bar S}$  is $\pi$-nef,
$-K_{\bar S}$  is $\q$-linearly equivalent to an effective linear combination of $\pi$-exceptional curves. \qed
\medskip

\begin{rem}  \label{antieff.lem.rem} This is especially strong if $S$ has quotient singularities. Then, by \cite[3.13]{ksb}, there is a unique, maximal
  $\pi':S'\to S$ such that $-K_{S'}$  is $\q$-linearly equivalent to an effective linear combination of $\pi'$-exceptional curves.
  This gives an effective procedure to construct all Q-modifications.
  \end{rem}


\begin{lem}\label{rigid.quot.mod.lem.1}
  Let $\pi:\bar S\to S$ be a Q-modification and  $\pi^m$ as in (\ref{P.res.ample.say}.1).
   Then   $\pi^m$ is a composite of a  
  sequence of blow ups
  $$
  \bar S^{\rm m}=:Y_r\to Y_{r-1}\to \cdots \to Y_1\to Y_0:= S^{\rm m},
  $$
  where each $Y_{i+1}\to Y_i$ is the blow-up of a singular point of the
  image if $\bar\Delta$ in $Y_i$.
\end{lem}

Proof. By (\ref{P.res.ample.say}.2)   $K_{\bar S^{\rm m}}\sim  \bar\mu^* K_{\bar S}-\bar\Delta$, 
 where
$\rdown{\bar\Delta}=0$ since $\bar S$ has only quotient singularities.

 For any $j$, let  $\bar\Delta_j$ denote the push-forward of $\bar\Delta$ to $Y_j$. Then   $K_{Y_r}+\bar\Delta_r\simq \bar\mu^* K_{\bar S}$ is nef over $S$
 by assumption, hence so is each
 $K_{Y_j}+\bar \Delta_j$.
 Let now $F_{i+1}$ denote the 
 exceptional curve of $Y_{i+1}\to Y_i$. Then $(K_{Y_{i+1}}\cdot F_{i+1})=-1$,
  so $(\bar \Delta_{i+1}\cdot F_{i+1})\geq 1$.  Since $\rdown{\bar\Delta_j}=0$, $F_{i+1}$ must intersect at least 2
  irreducible components of $\bar\Delta_{i+1}$ (different from $F_{i+1}$).
  The images of these in $Y_i$  show that the blow-up center is
  a singular point of $\bar\Delta_{i}$.
\qed





      

\medskip

These show that Q-modifications are combinatorial objects that are
determined by the dual graph ${\mathcal D}(S)$.
In fact, the following much stronger result holds, which generalizes  \cite[7.2]{MR1129040}.
To state it,
consider a graph with vertices $V$ and edges $E$. For a subset  $V'\subset V$ the  {\it induced 
  subgraph}  contains all the edges in $E$ that go  between vertices in $V'$.




  
\begin{prop}
  \label{Pmod.extends.prop}
  Let $S_1, S_2$ be surfaces with rational singularities
  and $j: {\mathcal D}(S_1)\into {\mathcal D}(S_2)$ an
isomorphism onto an induced 
subgraph that preserves the labeling of vertices.
Then there is a one-to-one correspondence between
    \begin{enumerate}
      \item Q-modifications $\tau_1: S_1^{\rm m}\to S_1\leftarrow  \bar S_1$ of $S_1$, and
      \item Q-modifications  $\tau_2: S_2^{\rm m}\to S_2\leftarrow \bar S_2$ of $S_2$,  such that $\ex(\tau_2)$ is  supported on
        $j( {\mathcal D}(S_1))$.
    \end{enumerate}
    Under this correspondence, $\bar S_1$ and $\bar S_2$ have the same singularities.
\end{prop}





Proof. Given a Q-modification  of $S_1$, we construct  a Q-modification  of $S_2$ as follows.

By (\ref{rigid.quot.mod.lem.1}),  $\pi_1^{\rm m}: \bar S_1^{\rm m}\to S_1^{\rm m}$ is obtained by repeatedly blowing up singular points of the exceptional divisor. We do the same blow ups to
get $\bar S_2^{\rm m}\to S_2^{\rm m}$. Then $j$ lifts to an
embedding $\bar{\jmath}$ of the dual graph of $\ex(\bar S_1^{\rm m}\to S_1)$
into the dual graph of $\ex(\bar S_2^{\rm m}\to S_2)$.

Next we get $\bar\mu_1: \bar S_1^{\rm m}\to \bar S_1$ 
by contracting certain
curves contained in
$\ex(\bar S_1^{\rm m}\to S_1)$.
Using $\bar{\jmath}$, we construct
$\bar\mu'_2: \bar S_2^{\rm m}\to \bar S'_2$
by contracting the  corresponding
curves of $\ex(\bar S_2^{\rm m}\to S_2)$.
By construction, $\bar S_2$ has the same singularities as $\bar S_1$.

We need  to show that $K_{\bar S_2}$ has nonnegative intersection number with
every exceptional curve of $\bar S_2\to S_2$.
For images of curves in $\ex(\pi_2^{\rm m})$ this follows from (\ref{P.res.ample.say}.5). For images of curves in $\ex(\mu_2)$ this follows from (\ref{P.res.ample.say}.4).

Conversely, let  $\tau_2: S_2^{\rm m}\to S_2\leftarrow \bar S_2$  be a  Q-modification of $S_2$ as in (\ref{Pmod.extends.prop}.2).
Then 
(\ref{rigid.quot.mod.lem.1}) shows that
$\bar S_2^{\rm m}\to S_2^{\rm m}$ is obtained by repeatedly blowing up
singular points of $j( {\mathcal D}(S_1))$. This shows that the above procedure can be reversed to get $\tau_1: S_1^{\rm m}\to S_1\leftarrow  \bar S_1$ of $S_1$.
\qed


\medskip






\begin{defn} \label{sim.Q.mod.defn} Let $p:X\to \dd$ be a flat morphism such that $X_0$ has rational singularities. A {\it simultaneous Q-modification} is a
  proper, birational morphism  $\pi:\bar X\to X$ such that
\begin{enumerate}
       \item $\pi_t:\bar X_t\to X_t$ is a Q-modification for every $t\in \dd$, and
    \item $K_{\bar X}$ is $\q$-Cartier. Equivalently,     $p\circ \pi:\bar X\to \dd$ is a KSB deformation.
\end{enumerate}
Being a Q-modification is an open property for KSB deformations.
Thus if  $\pi_0:\bar X_0\to X_0$ is a Q-modification, then the same holds for nearby $t$.
\end{defn}


The next result, proved in \cite[3.5]{ksb},  says that flat deformations of a surface  with quotient singularities always have a simultaneous Q-modification.
(See  \cite[5.41]{k-modbook} for higher dimensions.)

\begin{thm}\label{cmod.ksb.thm}
Let $p:X\to \dd$ be a flat morphism such that $X_0$ has quotient singularities. Then there is a unique  simultaneous Q-modification $\pi:\bar X\to X$ such that
  \begin{enumerate}
  \item $\pi$ is an isomorphism over $X\setminus X_0$, and 
        \item $K_{\bar X}$ is  $\pi$-ample.
    \item If $X_t$ has Du~Val singularities for $t\neq 0$, then
 $\pi_0:\bar X_0\to X_0$ is a P-modification  (\ref{P+M.defn}).
      \qed
  \end{enumerate}
\end{thm}






\section{P-modifications}

By \cite[3.9]{ksb}, for every quotient singularity $S$, the irreducible components of $\defsp(S)$ are in one-to-one correspondence with the 
P-modifications of $S$. We define and study these next.

\begin{defn}\label{P+M.defn}
  Let $(s,S)$ be a rational surface singularity. A proper, birational morphism $\pi_P:S_P\to S$ is a {\it P-modification}
if $S_P$ is normal,
\begin{enumerate}
\item $K_{S_P}$ is $\pi_P$-ample, and
\item $S_P$ has only Du~Val and $\a^2/\tfrac1{rn^2}(1, arn{-}1)$ singularities where  $(a,n)=1$.
\end{enumerate}
Following  \cite{be-ch}, $\pi_M:S_M\to S$ is an
 {\it M-modification} if  $S_M$ is normal,
\begin{enumerate}
\item $K_{S_M}$ is $\pi_M$-nef, and
\item $S_M$ has only  $\a^2/\tfrac1{n^2}(1, na{-}1)$ singularities, where  $n>1$ and $(a,n)=1$.
\end{enumerate}
Given a P-modification  $S_P\to S$, resolving its Du~Val  singularities and applying the construction (\ref{ksb.defs.say}) 
gives an M-modification $\pi_{MP}:S_M\to S_P$ such that
$K_{S_M}\simq \pi_{MP}^*K_{S_P}$. By \cite{be-ch} this establishes a
one to one correspondence between P- and M-modifications of $S$.

{\it Note on terminology.} For P-modifications we follow \cite[3.8]{ksb}.
A more general version is defined in \cite[6.2.13]{k-etc}, which allows more  singularities on $S_P$.
\end{defn}

For P-modifications  we get the following variant of (\ref{Pmod.extends.prop}).



\begin{cor}\label{Pmod.extends.prop.cor} Using the notation and assumptions of (\ref{Pmod.extends.prop}), there is a one-to-one correspondence between
    \begin{enumerate}
      \item P-modifications $\tau_1: S_1^{\rm m}\map  \bar S_1$ of $S_1$, and
      \item P-modifications  $\tau_2: S_2^{\rm m}\map \bar S_2$ of $S_2$,  such that $\ex(\tau_2)$ is  supported on
        $j( {\mathcal D}(S_1))\cup\{(-2)\mbox{ curves disjoint from }j( {\mathcal D}(S_1))\} $.
    \end{enumerate}
\end{cor}

Proof. The only difference is that, as we go from $\bar S_1$ to $\bar S_2$, we have to ensure that $K_{\bar S_2}$ is ample. That is, we have to contract
each connected component of  all $(-2)$-curves disjoint from $j( {\mathcal D}(S_1))$ to  a Du~Val singularity. \qed



\begin{say}[P-modifications of dihedral quotients]\label{dih.P.mod.all}  By \cite{MR1129040, MR1213355}
  P-modifications $\tau_D: S_D^{\rm m}\to S_D\leftarrow \bar S_D$ of $S_D:=S(2^2, {\mathbf c})$ are obtained from
  P-modifications  $\tau_C: S_C^{\rm m}\to S_C\leftarrow \bar S_C$ of  $S_C:=S(2, {\mathbf c})$ as follows.
    \begin{enumerate}
  \item ($C_1$ is not contracted by $\tau_C$.)  Then  $C'_0, C''_0$ are  contracted to  $A_1$ points by $\tau_D$.
  \item ($C'_0,C_1, \dots, C_i$ are contracted by $\tau_C$ to an $A_{i+1}$ point.)  Then  $C''_0$ is also  contracted  by $\tau_D$, giving a  $D_{i+2}$ point.
    \item ($C_1, \dots, C_i$ are contracted by $\tau_C$ to a non-DV point.)  Then  $C'_0, C''_0$ are  not contracted   by $\tau_D$.
    \item ($C'_0, C_1, \dots, C_i$ are contracted by $\tau_C$ to a non-DV point.)  Then  $C''_0$ is  not contracted  by $\tau_D$. There is also a symmetric version, where
      $C''_0,C_1, \dots, C_i$ are contracted by $\tau_D$ but $C'_0$ is  not contracted.
    \end{enumerate}
    \end{say}


  The P-modifications that appear in Theorem~\ref{dih.nplt.enum.i}  have several characterizations.


  \begin{prop} \label{DP.char.lem}
    Let $\pi_P:S_P\to S$ be a P-modification of a dihedral quotient singularity
    $(S^d_{n,q}, D^d_{n,q})$ with reduced exceptional curve $E_P$. The following are equivalent.
    \begin{enumerate}
    \item $S_P$ is one of the  cases (\ref{dih.P.mod.all}.1--2).
      \item $C'_0, C''_0$ are contracted to  Du~Val point(s) on $S_P$.
     \item   $(S_P, D_P+E_P)$ is log canonical. 
  \item   $K_{S_P}+D_P+E_P$ is numerically $\pi_P$-trivial.
    \end{enumerate}
\end{prop}

  Proof.  The equivalence of (1) and (2) is clear.
  To see (3) note that if $C_1, \dots, C_i$ are contracted by $\tau_C$ to a non-DV point, then there are 3 curves in $D_P+E_P$ through this point, and if
  $C'_0, C_1, \dots, C_i$ are contracted by $\tau_1$ to a non-DV point, then
  $C''_0$ meets the resolution at $C_1$, which is not  an end curve  of
  the chain for $i\geq 2$. For $i=1$ it is an  end curve,  but then another curve of $D_P+E_P$ meets $C_1$.

  For (4), note that $K_{S^{\rm m}}+\tsum_{i>0} C_i+\frac12 (C'_0+C''_0)$ is numerically trivial on $S^{\rm m}\to S$. Thus
  $K_{S_P}+D_P+E_P-\frac12 (C'_0+C''_0)$ is numerically $\pi_P$-trivial.
  Thus (4) holds iff both $C'_0,C''_0$ are contracted. This happens only in 
 cases (\ref{dih.P.mod.all}.1--2).\qed

  \begin{rem} \label{dih.P.mod.all.rem} Note that the cases (\ref{DP.char.lem}) are in bijection with
    P-modifications of  $T:=S(c_2,\dots, c_r)$.
    Indeed,  let $\tau_T:  T^{\rm m}\to T\leftarrow  T_P$
  be a P-modification.

    If $c_1>2$ then we get $\tau_D$ 
    by not contracting $C_1$ but contracting  $C'_0, C''_0$. 

    If $c_1=2$ and $C_2$ is contracted to a non-DV point by $\tau_T$, then again we contract  $C'_0, C''_0$ but not $C_1$. If $C_2,\dots, C_i$ are contracted to an $A_{i-1}$  point  by $\tau_T$, then $C'_0, C''_0, C_1, \dots, C_i$ are  contracted to a $D_{i+2}$ point  by $\tau_D$.
    Finally if  $C_2$ is not contracted, then 
    $C'_0, C''_0, C_1$  are  contracted to an $A_{3}$ point  by $\tau_D$.
  \end{rem}





The main  reason to study 
P- and M-modifications is the following, essentially proved in \cite[3.9]{ksb} and  \cite{be-ch}.

\begin{thm}\label{main.MP.old.thm} Let $(0,S)$ be a  rational surface singularity,  $\pi_{P}:S^{}_P\to S$  a P-modification,  and
  $\pi_{MP}:S^{}_M\to S^{}_P$ the corresponding M-modification.
  Then there is a commutative diagram
  $$
\begin{array}{ccccc}
\univ_{\rm KSB} (S^{}_M) & \stackrel{c^{}_{MP}}{\longrightarrow} &
  \univ_{\rm KSB} (S^{}_P) & \stackrel{c^{}_{P}}{\longrightarrow} & \univ (S)\\[1ex]
 u^{}_M \downarrow\hphantom{u^{}_M}  && u^{}_P \downarrow\hphantom{u^{}_P}  && u_S \downarrow\hphantom{u_S} \\
  \defsp_{\rm KSB} (S^{}_M) & \stackrel{\tau^{}_{MP}}{\longrightarrow} &
  \defsp_{\rm KSB} (S^{}_P) & \stackrel{\tau^{}_{P}}{\longrightarrow} &
  \defsp (S)
\end{array}
\eqno{(\ref{main.MP.old.thm}.1)}
$$
where
\begin{enumerate}\setcounter{enumi}{1}
  \item $\defsp_{\rm KSB}(S^{}_M)$ and
    $\defsp_{\rm KSB}(S^{}_P)$ are smooth of the same dimension,
  \item   $\tau^{}_{MP}: \defsp_{\rm KSB}(S^{}_M) \to \defsp_{\rm KSB}(S^{}_P)$ is finite and Galois,
  \item $\tau^{}_{P}: \defsp_{\rm KSB}(S^{}_P) \to \defsp(S)$
    is a finite, birational morphism onto an irreducible component of
    $\defsp(S)$.
    \end{enumerate}
\end{thm}


Sketch of proof. For the existence of the diagram we need to show that the
morphisms  $\pi_{P}:S^{}_P\to S$   and
$\pi_{MP}:S^{}_M\to S^{}_P$ extend to
their flat deformations
$$
c^{}_{P}:  \univ_{\rm KSB} (S^{}_P)\to \univ (S)
\qtq{and}
c^{}_{MP}: \univ_{\rm KSB} (S^{}_M) \to \univ_{\rm KSB} (S^{}_P).
$$
The relevant property is the vanishing of
$R^1(\pi_{P})_*\o_{S^{}_P}$ and of $R^1(\pi_{MP})_*\o_{S^{}_M}$. 
The proof is summarized in \cite[11.4]{km-flips}, which in turn relies mainly
on \cite{MR54:10261}. 


Since KSB deformations are locally unobstructed, the obstruction to smoothness lies in $H^2$ of the tangent sheaf. If $X\to U$ is proper of fiber dimension $\leq 1$, and $U$ is affine, then $H^i$ of any coherent  sheaf on $X$ vanishes for $i\geq 2$. This implies smoothness in (\ref{main.MP.old.thm}.2); see (\ref{loc.glob.lem}) for details.

The relation between
P- and M-modifications is established in \cite{be-ch}.
The Galois group in  (\ref{main.MP.old.thm}.3) is a product of reflection groups, determined by the singularities of $S^{}_P$, see \cite{be-ch} and (\ref{ksb.defs.say}.6) for a key example.



In order to show (\ref{main.MP.old.thm}.4), note first that 
every fiber of $u^{}_P$ is a P-modification of the corresponding deformation of $S$. We check in   (\ref{rigid.quot.mod.thm}) that P-modifications $\pi_{P}:S^{}_P\to S$ do not have nontrivial 1-parameter deformations fixing $S$.
This shows that  $\tau^{}_{P}$ is finite.
As in \cite{art-bri},  $\tau^{}_{P}$
    is  onto an irreducible component of $\defsp(S)$
    by openness of versality.
Thus $\tau^{}_{P}: \defsp_{\rm KSB}(S^{}_P) \to \defsp^{}(S)$
is a finite cover. In order to show that it has degree 1, let
$T\to \defsp^{}(S)$ be a morphism from the spectrum of a DVR that maps the closed (resp.\ generic) point  the closed (resp.\ generic) point.
By what we already proved, there is a finite $T'\to T$ such that the pull-back of $\univ(S)$ to $T'\to T\to \defsp^{}(S)$ is obtained from a KSB deformation of  $S^{}_P$. We may assume that $T'/T$ is Galois. As noted in \cite[p.312]{ksb}, 
the action of $\gal(T'/T)$ lifts to an action on this KSB deformation.
Taking the quotient shows that the pull-back of $\univ(S)$ to $T\to \defsp^{}(S)$ is obtained from a KSB deformation of  $S^{}_P$. Thus
$T\to \defsp^{}(S)$ lifts to $T\to \defsp_{\rm KSB}(S^{})$, so
  $\tau^{}_{P}$ is birational. \qed
    
    
    
\medskip
{\it Remark \ref{main.MP.old.thm}.5.} By  \cite{MR1129040},  $\tau^{}_{P}$
is  an isomorphism onto an irreducible component of $\defsp(S)$ for cyclic quotients, but
there are dihedral examples where $\tau^{}_{P}$
is  not an isomorphism.
Also, $\tau^{}_{P}$ is birational but not necessarily finite
 for the more general 
P-modifications considered in \cite[Sec.6]{k-etc}.



\section{Tangent sheaves}


In order to compute the dimensions of the deformation spaces, 
we need various results on  tangent and logarithmic tangent sheaves for
modifications of rational singularities. These are mostly taken from
\cite{MR367277, MR333238} and 
\cite{lee-park, MR2995167}.  Note that \cite{lee-park, MR2995167} focus on the global aspects, which tend to be more subtle. 
See also \cite{pa-pa-sh-2, MR4509033, MR4502069} for similar computations. 


\begin{say}[Tangent and logarithmic tangent sheaves]\label{log.tg.sh.say}
  Let  $S$ be a smooth surface and $D\subset S$ a smooth curve with normal bundle $N_D$.
  There is a natural map $T_S\to N_D$. Its kernel is the
  {\it logarithmic tangent bundle} of the pair $(S, D)$, denoted by
  $T_S(-\log D)$.  Thus we have
  an exact sequence
  $$
  0\to T_S(-\log D) \to T_S\to N_{D}\to 0.
  \eqno{(\ref{log.tg.sh.say}.1)}
  $$
  The restriction of $  T_S(-\log D)$ to $D$ sits in an exact sequence
  $$
  0\to \o_D\to T_S(-\log D)|_D \to T_D\to 0.
  \eqno{(\ref{log.tg.sh.say}.2)}
  $$
  Locally, if $D=(y=0)$, then the kernel  $\o_D$ is generated by
  $y\tfrac{\partial}{\partial y}$ and the quotient $T_D$ by
  $\tfrac{\partial}{\partial x}$.

  Next assume that  $S$ is smooth but $D\subset S$ is a nodal curve
  with normalization $\tau: \bar D\to D$. The {\it immersed normal sheaf} of $D$ is
  $$
  N_{\bar D}:=  \tau_*\coker\bigl[T_{\bar D}\to \tau^*T_S\bigr].
   \eqno{(\ref{log.tg.sh.say}.3)}
 $$
 Then (\ref{log.tg.sh.say}.1) becomes
  $$
  0\to T_S(-\log D) \to T_S\to N_{\bar D}\to 0,
 %%  \eqno{(\ref{log.tg.sh.say}.4)}
 $$
   where $T_S(-\log D)$ is locally free. More generally, if $D=D_1+D_2$, then we have
$$
  0\to T_S(-\log (D_1+D_2)) \to T_S(-\log D_1)\to N_{\bar D_2}\to 0.
   \eqno{(\ref{log.tg.sh.say}.4)}
 $$
   The sequence (\ref{log.tg.sh.say}.2) now gives
   $$
  0\to \o_{\bar D}\to \tau^*T_S(-\log D) \to T_{\bar D}(-N)\to 0,
  \eqno{(\ref{log.tg.sh.say}.5)}
  $$
  where $N\subset \bar D$ is the preimage of the nodes of $D$.
  
  If $S$ is a normal surface and $D\subset S$ a reduced curve, then there is a finite subset $P\subset S$ such that $(S\setminus P, D\setminus P)$ is a smooth pair. Let $j:S\setminus P\into S$ be the natural injection.
  The {\it tangent sheaf} of $S$  and the {\it logarithmic tangent sheaf} of $(S, D)$ are defined as
  $$
  T_S:=j_*T_{S\setminus P}, \qtq{and}
  T_S(-\log D):= j_*\bigl(T_{S\setminus P}(-\log (D\setminus P))\bigr).
   \eqno{(\ref{log.tg.sh.say}.6)}
   $$
   In the log canonical case, we have the following close analog of (\ref{log.tg.sh.say}.4).

  
\end{say}

\begin{lem} \label{log.tg.sh.thm}
  Let $(S, D_1+D_2)$ be a log canonical pair, 
   $\mu:S^{\rm m}\to S$  the minimal resolution, and $D^{\rm m}_i:=\mu^{-1}_*D_i$ with 
normalizations  $\tau_i: \bar D^{\rm m}_i\to D^{\rm m}_i$. Then 
there is an exact sequence
  $$
0\to T_S(-\log (D_1+D_2)) \to T_S(-\log (D_1))\to \mu_*N_{\bar D^{\rm m}_2}\to 0.
 \eqno{(\ref{log.tg.sh.thm}.1)}
  $$
\end{lem}

Proof.  Over the smooth locus of $S$, this is (\ref{log.tg.sh.say}.4).
Thus it remains to check what happens at the singular points that are contained in $D_1+D_2$. For this we can use local analytic coordinates.


  If a finite group $G$ acts with isolated fixed points on a surface $X$, then the local sections of $T_{X/G}$ are the $G$-invariant local sections of $T_X$.
Thus the tangent sheaf of $S_{n,q}=\a^2_{xy}/\frac1{n}(1,q)$ is generated by
  $$
  x\tfrac{\partial}{\partial x}, y^{q'}\tfrac{\partial}{\partial x}, y\tfrac{\partial}{\partial y}, x^q\tfrac{\partial}{\partial y}.
   \eqno{(\ref{log.tg.sh.thm}.2)}
   $$
   If $B_x\subset S_{n,q}$ is the image of the $x$-axis, then we see that
   $T_S(-\log B_x)$ is generated by
   $x\tfrac{\partial}{\partial x}, y^{q'}\tfrac{\partial}{\partial x},  y\tfrac{\partial}{\partial y}$, and the quotient
   $T_S/T_S(-\log B_x)$ is generated by $x^q\tfrac{\partial}{\partial y}$.
   

   For the minimal resolution of $S_{n,q}$ we use  the chart given in
   (\ref{res.of.cyclic.say}.2). The chain rule  gives that
   $$
   (\pi_1)_* \bigl(\tfrac{\partial}{\partial y_2}\bigr)=x^q\tfrac{\partial}{\partial y}.
   $$
   That is, $T_S/T_S(-\log B_x)$ is naturally isomorphic to the
   normal bundle of $B_x^{\rm m}\subset S^{\rm m}_{n,q}$. This shows (\ref{log.tg.sh.thm}.1) at the cyclic quotient points.

   At a dihedral point  $(S^d_{n,q}, B^d_{n,q})$, the  minimal resolution can be obtained by first taking a double cover
   $(S_{N,Q}, D_{N,Q})\to (S^d_{n,q}, B^d_{n,q})$, resolving $S_{N,Q}$ and then quotienting out by the involution. The latter is fixedpoint-free along the birational transform of $D_{N,Q}$, so the normal bundle computation for the birational transform of $B^d_{n,q}$ is the same as for the birational transform of $D_{N,Q}$.
   Finally note that $Q^2\equiv 1\mod N$, so  the 
   quotient  $T_S/T_S(-\log B_x)$ is generated by 
   $x^Q\tfrac{\partial}{\partial y}+y^Q\tfrac{\partial}{\partial x}$,
   where $x,y$ are the  orbifold  coordinates on $S_{N,Q}$.
   This shows (\ref{log.tg.sh.thm}.1) at the dihedral quotient points. \qed
   
   

   \begin{prop}\label{cyc.lc.lf.prop}
     Let $(X, D)$ be a log canonical pair with  quotient singularities.  Let $\pi:X^{\rm m}\to X$ be either the minimal resolution,
or the blow-up of a node of $D$ where $X$ is smooth.
    Let   $E^{\rm m}$ be the exceptional curve and  $D^{\rm m}$ the birational transform of $D$. Then
     \begin{enumerate}
     \item $\pi_*T_{X^{\rm m}}\bigl(-\log (E^{\rm m}+D^{\rm m})\bigr)=T_X(-\log D)$,\item $R^1\pi_*T_{X^{\rm m}}\bigl(-\log (E^{\rm m}+D^{\rm m})\bigr)=R^1\pi_*T_{X^{\rm m}}\bigl(-\log (E^{\rm m})\bigr)= 0$, and
     \item  $H^i\bigl(X, T_X(-\log D)\bigr)=
    H^i\bigl(X^{\rm m}, T_{X^{\rm m}}\bigl(-\log (E^{\rm m}+D^{\rm m})\bigr)\bigr)$ for every $i$.
     \end{enumerate}
   \end{prop}

  
  

   
 Proof.
 For (\ref{cyc.lc.lf.prop}.1) the 
 blow-up of a node of $D$ is a simple computation.
At a  singular point of $X$, 
 it suffices to show the $D=\emptyset$ case. That is,  all local generators of
   $T_X$  lift to sections of
 $T_{X^{\rm m}}\bigl(-\log E^{\rm m}\bigr)$.  
 This should hold for all normal singularities (\ref{TX.list.rem}), but here is an argument for quotient  singularities.


 Write $S=\a^2/G$. Let $\pi:B_0\a^2\to \a^2$ the blow-up of the origin with  exceptional curve $E$. By direct computation,
 $\pi_*T_{B_0\a^2}(-\log E)= m_0T_{\a^2}$,
 where $m_0$ is the maximal ideal at the origin.
Set $\pi':S':=B_0\a^2/G\to S$ with  exceptional curve $E'$.
 Taking $G$-invariants gives that
$\pi_*T_{S'}(-\log E')= (m_0T_{\a^2})^G$.
 Now note that the $G$-action on $T_{\a^2}/m_0T_{\a^2}$ has no fixed vectors,
 so $(m_0T_{\a^2})^G=T_{\a^2}^G=T_S$.

 The $G$-action on $B_0\a^2$ leaves $E$ invariant, so 
 $S'$ has only cyclic quotient singularities.
 Unfortunately, iterating  the above process does not give a resolution, see \cite[2.29]{k-res}. However, for cyclic quotients we can use the blow-up charts
(\ref{res.of.cyclic.say}.1--2).  By  direct computation
   $
   x\tfrac{\partial}{\partial x}=
   nx_2\tfrac{\partial}{\partial x_2}-q y_2\tfrac{\partial}{\partial y_2},
   $
   proving (\ref{cyc.lc.lf.prop}.1).

The $D=\emptyset$ case of 
(\ref{cyc.lc.lf.prop}.2)  is discussed in (\ref{taut.defn}.1).
For general $D$, we use (\ref{log.tg.sh.say}.4) to get
$$
0\to  T_{X^{\rm m}}\bigl(-\log (E^{\rm m}+D^{\rm m})\bigr)
\to  T_{X^{\rm m}}\bigl(-\log (E^{\rm m})\bigr)
\to  N_{D^{\rm m}}\to 0.
$$
Pushing this forward, we get that
$$
R^1\pi_*T_{X^{\rm m}}\bigl(-\log (E^{\rm m}+D^{\rm m})\bigr)\cong
\coker\bigl[T_S\to \pi_*N_{\bar D^{\rm m}}\bigr],
$$
  and the latter is 0 by (\ref{log.tg.sh.thm}).
   The Leray spectral sequence now gives (\ref{cyc.lc.lf.prop}.3).\qed



   \medskip
       {\it Remark \ref{cyc.lc.lf.prop}.4.}   Let $\pi:X'\to X$ be the blow-up of a smooth point with exceptional curve $E$. Then
       $h^1(X', T_{X'}(-\log E))=h^1(X, T_X)+2$.
       
   \medskip
{\it Complement \ref{cyc.lc.lf.prop}.5.}  The computations also show that if all the singularities are 
$(S_{n,q}, D_{n,q})$, then  $T_X(-\log D)$ is locally free
and $\pi^*T_X(-\log D)\cong  T_{X^{\rm m}}\bigl(-\log (E^{\rm m}+D^{\rm m})\bigr)$.



\begin{ques} \label{TX.list.rem}
  Let $X$ be a normal variety and  $\pi:X^{\rm c}\to X$ its canonical modification.  Every automorphism of $X$ lifts to an automorphism of $X^{\rm c}$, and, if $X$ is proper, then  $H^0(X, T_X)$ is the tangent space of the automorphism group.  Thus one can expect that
  $$
  \pi_*T_{X^{\rm c}}=T_X.
  \eqno{(\ref{TX.list.rem}.1)}
  $$
  \cite[4.1]{MR1320605} says that if $X$ is local and complete, then again
  $H^0(X, T_X)$ is the tangent space of the automorphism group, suggesting that 
  (\ref{TX.list.rem}.1) indeed holds.

  The caveat is that  both objects are infinite dimensional, so one would need to check that the above finite-dimensional intuition  works out in general. 
  \end{ques}



We can now compute the global contribution to the dimension of  our
deformation spaces.



\begin{prop}\label{dim.h1T.prop.ol}
  Let $S$ be a affine  surface with rational singularities and $D\subset S$ a reduced curve. Let 
  $\mu:S^{\rm m}\to S$ be the  minimal resolution with  reduced exceptional
  set $E^{\rm m}$, and 
  $\pi:{\bar S}\to S$  a Q-modification with  reduced exceptional set $\bar E$. Let  $D^{\rm m}$ resp.\ $\bar D$ denote the birational transforms of $D$ on $S^{\rm m}$  resp.\ $\bar S$. Assume that $(\bar S, \bar E+\bar D)$ is lc and
  $(S^{\rm m},E^{\rm m}+D^{\rm m})$ is normal crossing.
  
For an irreducible component $\bar E_i\subset \bar E$, let
  $e_i\in \n$ be the negative of the self-intersection of its
birational transform on the minimal resolution of ${\bar S}$.
Then
\begin{enumerate}
\item $h^1\bigl({\bar S}, T_{\bar S}(-\log \bar D)\bigr)=
  h^1\bigl({\bar S}, T_{\bar S}(-\log ( \bar E+ \bar D))\bigr)+
  \tsum_{i\in I} (e_i-1)$, and
  \item 
    $h^1\bigl({\bar S}, T_{\bar S}(-\log ( \bar E+ \bar D))\bigr)=
    h^1\bigl(S^{\rm m}, T_{S^{\rm m}}(-\log (E^{\rm m}+D^{\rm m}))\bigr).$
\end{enumerate}
\end{prop}




Proof.  
(\ref{log.tg.sh.thm}) gives an exact sequence
$$
0\to  T_{\bar S}\bigl(-\log ( \bar E+ \bar D)\bigr)\to
T_{\bar S}(-\log \bar D)\to \oplus \o_{\bar E_i}(-e_i)\to 0.
$$
Here $h^0\bigl(\bar E_i, \o_{\bar E_i}(-e_i)\bigr)=0$ since $e_i>0$, so taking cohomologies gives the first claim.
For the second, we use the notation of (\ref{P.res.ample.say}), and claim that
$$
h^1\bigl(\bar S^{\rm m}, T_{\bar S^{\rm m}}(-\log (\bar F^{\rm m}+\bar D^{\rm m}))\bigr)
$$
equals  both terms in (\ref{dim.h1T.prop.ol}.2), where $\bar F^{\rm m} $ is the reduced exceptional set of  $\bar S^{\rm m}\to   S$.
For $\bar S^{\rm m}\to  \bar S$ we use the minimal resolution case of (\ref{cyc.lc.lf.prop}).
Next, since  $\bar S^{\rm m}\to  S^{\rm m}$ is a composite of node blow-ups by (\ref{rigid.quot.mod.lem.1}), we can repeatedly use the
node blow-up case of (\ref{cyc.lc.lf.prop}).  \qed



\section{Dimension formulas}


We compute the dimension of the irreducible components of various deformation spaces using invariants of the singularity and its P-modifications.
The formulas are simplest when $S$ is determined by ${\mathcal D}(S)$. 


\begin{say}
  \label{taut.defn}
Let  $(s, S)$ be a rational surface singularity with  
 minimal resolution  $\mu:S^{\rm m}\to S$ and  reduced exceptional curve
$C^{\rm m}=\cup_{i\in I} C_i$.

By \cite[Sec.3.4.4]{sernesi}, $H^1\bigl(S^{\rm m}, T_{S^{\rm m}}(-\log C^{\rm m})\bigr)$
is the tangent space to those deformations of $S^{\rm m}$ where every $C_i$ lifts. Equivalently, these are  those deformations of $S$ that preserve the dual graph.  The dimension of this tangent space is usually hard to compute, but it is known in several important cases.

\medskip
{\it Quotient singularities \ref{taut.defn}.1.} The dual graph determines
the singularity, hence $H^1\bigl(S^{\rm m}, T_{S^{\rm m}}(-\log C^{\rm m})\bigr)=0$.
This has been long known; see for example \cite{briesk},
  \cite[III.5.1]{bpv} or 
  \cite[3.32]{kk-singbook}.

  \medskip
      {\it Taut singularities \ref{taut.defn}.2.} These are the singularities that are determined by their dual graphs.
      For these $H^1\bigl(S^{\rm m}, T_{S^{\rm m}}(-\log C^{\rm m})\bigr)=0$.
      A rather laborious result of \cite{MR367277, MR333238} says that
 this vanishing implies tautness,  with a few exceptions; these are enumerated in \cite[3.2]{MR333238}.
 The quite long  complete list of taut singularities is given in \cite{MR333238}.

 A simple case is when  ${\mathcal D}(S)$ has only 1 fork, which has  degree 3 and self-intersection $\leq -3$. This is  used for the  $W(p,q,r)$ singularities studied in \cite{park2022deformations}.

   \medskip
       {\it Weighted homogeneous singularities \ref{taut.defn}.3.}
       For these  ${\mathcal D}(S)$ has only 1 fork. Let $d_0$ denote its   degree  and  $-c_0$ its self-intersection.  By \cite[4.1.III]{MR333238}, if
       $c_0\geq 2d_0-3$ then
       $H^1\bigl(S^{\rm m}, T_{S^{\rm m}}(-\log C^{\rm m})\bigr)=d_0-3$.

        These are used
        in \cite{jeon2023deformations} under the (mostly) weaker assumption
        $c_0\geq d_0+3$.
\end{say}

For  P-modifications, we need only the simplest numerical 
invariants.




\begin{defn}\label{I.J.P.mod.defn}   Let  $\pi:S_P\to S$ be a
  P-modification of a normal singularity. Let
$$
\{p_j\in S_P: j\in J_P\} \qtq{resp.} \{E_i\subset S_P: i\in I_P\}
$$
 denote the singular points of $S_P$, respectively
 the irreducible, exceptional curves of $\pi:S_P\to S$.
 Set $r_j=r$ if $p_j$ is of type $A_r, D_r, E_r$  or 
 $\a^2/\tfrac1{rn^2}(1, arn{-}1)$ for some $n>1$.
Note that $r_j=\dim \defsp_{\rm KSB}(p_j,S_P)$ by (\ref{ksb.defs.say}.2).
 
Let $e_i\in \n$ denote the negative of the self-intersection of the
birational transform of $E_i$ on the minimal resolution of $S_P$.
\end{defn}

We can now state the first dimension formula.




\begin{prop}\label{taut.dim.thm} Let $\pi:S_P\to S$ be a  P-modification of an affine surface  with  rational  singularities, and  $I_P, J_P$  as in (\ref{I.J.P.mod.defn}).
  Let $S^{\rm m}\to S$ be the minimal resolution, with reduced exceptional curve $C^{\rm m}$. 
  Then 
  $$
  \dim \defsp_{\rm KSB}(S_P)=\tsum_{j\in J_P}r_j+\tsum_{i\in I_P} (e_i-1)+h^1\bigl(S^{\rm m}, T_{S^{\rm m}}(-\log C^{\rm m})\bigr).
  \eqno{(\ref{taut.dim.thm}.1)}
  $$
  In particular, if $S$ has  taut singularities, then
  $$
  \dim \defsp_{\rm KSB}(S_P)=\tsum_{j\in J_P}r_j+\tsum_{i\in I_P} (e_i-1).
  \eqno{(\ref{taut.dim.thm}.2)}
  $$
\end{prop}






Proof.   First note that, by (\ref{loc.glob.lem}),
$$
\dim \defsp_{\rm KSB}(S_P)= \tsum_{j\in J_P} \dim  \defsp_{\rm KSB}(p_j,S_P)
+ h^1(S_P, T_{S_P}).
$$
Here  $r_j=\dim \defsp_{\rm KSB}(p_j,S_P)$ by (\ref{ksb.defs.say}.2),
this explains the summand $\tsum_{j\in J_P}r_j$. 
 The computation of $ h^1(S_P, T_{S_P})$ is the  $D=0$ special case of  (\ref{dim.h1T.prop.ol}). \qed





\medskip
For KSB deformations of pairs we have the following.
The dihedral case is discussed in (\ref{main.cyclic.thm.d}).
 

   \begin{thm}\label{main.cyclic.thm} Let $\pi:(S_P, D_P+E_P)\to (S,D):=(S_{n,q}, D_{n.q})$ be a P-modification of a cyclic quotient  singularity and 
     $J_P$  as  in (\ref{I.J.P.mod.defn}).
     Then
     $$
     \dim \defsp_{\rm KSB}(S_P, D_P+E_P)= \tsum_{j\in J_P}r_j+
     |\mbox{type $A$  points}|.
     \eqno{(\ref{main.cyclic.thm}.1)}
     $$
   \end{thm}

   {\it Clarification \ref{main.cyclic.thm}.2.} Here
   $ |\mbox{type $A$  points}|$ is the number of  points 
    where the pair $(S_P, D_P+E_P)$ is locally of the form
   $\bigl((xy=z^{r+1}), (xy=0)\bigr)$ for some $r\geq 0$.
   For $r=0$ the surface is smooth but the curve is singular. 
   \medskip

   Proof. If $p_j$ is a singular point of $S_P$ or of $D_P+E_P$, then
   $\dim \defsp_{\rm KSB}(p_j, S_P, D_P+E_P)$  is  $r$ for
   a singularity is of type  $\a^2/\tfrac1{rn^2}(1, arn{-}1)$ by (\ref{ksb.defs.pair.dd}.2), but
   $r+1$ if the singularity is of type  $A_r$.  (Here we need to count
   $\bigl(\a^2, (xy=0)\bigr)$ as type $A_0$.) Thus the
   right hand side of (\ref{main.cyclic.thm}.1) is the sum of the local terms in (\ref{loc.glob.lem}).

   It remains to show that the global term
   $h^1\bigl(S_P, T_{S_P}(-\log (D_P+E_P))\bigr)$ vanishes.
   This follows from (\ref{dim.h1T.prop.ol}.2),  (\ref{cyc.lc.lf.prop}.2) and
   (\ref{taut.defn}.1.). \qed



   \begin{rem} \label{dim.drop.cyc.say} Comparing (\ref{taut.dim.thm}) and
     (\ref{main.cyclic.thm}) we see that
     $$
     \dim \defsp_{\rm KSB}(S_P, D_P+E_P)\leq \dim \defsp_{\rm KSB}(S_P),
     $$
     and the inequality is strict with a few exceptions. These are
     triple points  (\ref{KSB.codin.infini.thm}) and also P-resolutions of the form
   $$
    \bullet {-\boxed{A_{n_1}}-}\ 2\ {-\boxed{\tfrac1{rn^2}(1, arn{-}1)}-} \ 2\
     {-\boxed{A_{n_2}}-} \bullet
     $$
     \end{rem}
   
  
The dihedral version of (\ref{main.cyclic.thm}) is the following.
  
   \begin{thm}\label{main.cyclic.thm.d} Let $\pi:(S_P, D_P+E_P)\to (S,D)=(S^d_{n,q}, D^d_{n.q})$ be a P-modification of a dihedral quotient  singularity as in
     (\ref{DP.char.lem}) and 
     $J_P$  as  in (\ref{I.J.P.mod.defn}).
     Then
     $$
     \dim \defsp_{\rm KSB}(S_P, D_P+E_P)= \tsum_{j\in J_P}r_j+
     |\mbox{type $A$  points}| -2.
     \eqno{(\ref{main.cyclic.thm.d}.1)}
     $$
   \end{thm}

  

   Proof.  Note that here $(S_P, D_P+E_P)$ has a singular point of type
   $$
   \bigl((z^2=x(y^2-x^{r-2})), (x=z=0)\bigr)$$  
   whose KSB deformation space has dimension $r{-}2$ by  (\ref{D.term.QC.exmp.1}).  (We have 2 singular points if $r=2$.)
   This is why we need to subtract 2 on
the
   right hand side of (\ref{main.cyclic.thm.d}.1). The rest of the argument is as for (\ref{main.cyclic.thm}). \qed


   

\medskip
We used the  following result, whose
 proof  is summarized in \cite[11.4]{km-flips}.

\begin{lem}\label{loc.glob.lem}
  Let $S$ be an affine  surface,  
  $\tau:X\to S$ a proper, birational morphism and $D\subset X$ a reduced curve. Assume that $X$ is normal and let $\{p_j: j\in J\}$ be the   points  of $\sing X\cup \sing D$. Then the restriction maps
  $$
  \defsp (X,D)\to \times_{j\in  J} \defsp (p_j, X, D)
  \qtq{and}
  \defsp_{\rm KSB} (X,D)\to \times_{j\in  J} \defsp_{\rm KSB}  (p_j, X,D)
  $$
 are smooth of relative dimension $h^1\bigl(X, T_X(-\log D)\bigr)$.  \qed
\end{lem}
  


By (\ref{ord.of.K+C.say}.1), the pairs  $(S_{n,q}, B_{n,q})$ are KSB rigid. By contrast, we show that the pairs  $(S_{n,q}, D_{n,q})$ are KSB smoothable, and the pairs 
$(S^d_{n,q}, D^d_{n,q})$ have KSB deformations whose general fibers have only
singularities of type 
$(S_{2,1}, D_{2,1})\cong  \bigl((xy=z^2), (x=z=0)\bigr)$; we call these
{\it almost smoothings.} More precisely, we have the following.

\begin{prop} \label{smoothing.prop}
  Every irreducible component of
  $\defsp_{\rm KSB} (S_{n,q}, D_{n,q})$ contains smoothings, and 
every irreducible component of
   $\defsp_{\rm KSB} (S^d_{n,q}, D^d_{n,q})$ contains almost smoothings.
\end{prop}



Proof.   By (\ref{ord.of.K+C.say}.4--5),  $K_S+D$ (resp.\ $2(K_S+D)$)  is Cartier, so the same holds for its KSB deformations, hence every singularity of a KSB deformation of our cyclic (resp.\ dihedral) pair is again   cyclic (resp.\ dihedral). By openness of versality
\cite{art-bri}, it is thus enough to show that each  cyclic (resp.\ dihedral) pair is smoothable  (resp.\ almost smoothable).

We find such deformations in the Artin component.
For a cyclic pair $(S, D)$, 
 look at the minimal resolution  
  $\pi_0:(S^{\rm m}, D^{\rm m}+E^{\rm m})\to (S, D)$. By (\ref{loc.glob.lem}),
 there is a deformation  $p: ({\mathbf S}^{\rm m}, {\mathbf D}^{\rm m})\to \dd$
 whose special fiber is $(S^{\rm m}, D^{\rm m}+E^{\rm m})$ and whose generic fiber is a smooth pair with ${\mathbf D}^{\rm m}_{gen}$ irreducible.
 Next $\pi_0$ extends to a contraction
 $$
 \pi:({\mathbf S}^{\rm m}, {\mathbf D}^{\rm m})\to
 ({\mathbf S}, {\mathbf D})\to \dd,
 $$
and $({\mathbf S}, {\mathbf D})\to \dd$ is a KSB smoothing
 since $K_{S^{\rm m}}+ D^{\rm m}+E^{\rm m}\sim 0$  by (\ref{ord.of.K+C.say}.4).


 For a dihedral pair $(S, D)$,  we start with
 $\pi_0:(S', D'+E')\to (S, D)$ which is obtained from the
 minimal resolution by contracting the curves  $C'_0, C''_0$  (as in Notation~\ref{dih.basic.not}).
  The key point is that $2\bigl(K_{S'}+ D'+E'\bigr)\sim 0$
  by (\ref{ord.of.K+C.say}.5).
  The above argument now works for this case too, except that we have 2 singularities of type
 $\bigl((xy=z^2), (x=z=0)\bigr)$. These are rigid, so persist in every deformation. 
\qed












\section{Doubly KSB deformations}\label{0-1.sec}




\begin{defn}\label{2-KSB.defn} A flat deformation 
  $p: {\mathbf S}\to \dd$   of $S\cong {\mathbf S}_0$  is a {\it doubly KSB deformation} of $(S, D)$ if
  it is a KSB deformation both for $S$ and for  $(S, D)$.
  Equivalently, both $K_S$ and $D$ lift to $\q$-Cartier $\z$-divisors on
  ${\mathbf S}$.

  Since we have an explicit description of all  KSB deformations,
  we can use it to describe all doubly KSB deformations in most cases. The basic examples are the following.
\begin{enumerate}
 \item Every KSB smoothing of $S_{n,q}$ is a
  doubly KSB deformation of  $(S_{n,q}, D_{n,q})$ by (\ref{ksb.defs.pair.dd}.2).
\item By Proposition~\ref{cyc.plt.rig.i}, $(S_{n,q}, B_{n,q})$ has no doubly KSB deformations.
\item 
  If  $\frac56<d\leq 1$ and $(S,dD)$ is lc, then $(S, D)$ is also lc, see \cite[3.44]{kk-singbook}.
  Thus the doubly KSB smoothings are fully described by (\ref{2-KSB.defn}.1--2).
\item Other examples with $d=\frac12$ are in (\ref{1/2.2.exmp});
  I did not find any  with $\frac12< d<1$. See also (\ref{D.term.QC.exmp.1}).
\end{enumerate}

\end{defn}




\begin{thm}\label{K+D.KSB.P.thm}  Let $(S, dD)$ be an lc pair such that
  either $d>\frac12 $ or $d=\frac12$ and   $(S, dD)$ is klt.
  There is a bijection between the following sets.
\begin{enumerate}
\item  \label{K+D.KSB.P.thm.1}
  Irreducible components of $\defsp_{\rm KSB}(S, dD)$ whose generic fiber has only Du~Val singularities.
\item \label{K+D.KSB.P.thm.2}
  P-modifications $\pi_P:S_P\to S$ such that
  \begin{enumerate}
\item  \label{K+D.KSB.P.thm.2.a} $K_{S_P}+d(D_P+E_P)\simq 0$, and
\item   \label{K+D.KSB.P.thm.2.b} the singularities of $\bigl(S_P, D_P+E_P\bigr)$ have doubly KSB deformations with Du~Val  generic fibers.
 \end{enumerate}
\end{enumerate}
\end{thm}

Proof.   Assume that  $\pi_P:S_P\to S$ satisfies (\ref{K+D.KSB.P.thm}.2).
By (\ref{loc.glob.lem})  the local doubly KSB deformations  glue together to a doubly KSB deformation of
$\bigl(S_P, d(D_P+E_P)\bigr)$. The contraction $\pi_P$ extends to nearby deformations as in (\ref{main.MP.old.thm}), and we get a whole irreducible component by  openness of versality \cite{art-bri}.  (Note that this part holds without any restriction on $d$.)

Conversely, take a general 1-parameter KSB deformation
$p: ({\mathbf S}, d{\mathbf D})\to \dd$ of $(S, dD)$.
Let   $\pi: \bar {\mathbf S}\to {\mathbf S}$ be the simultaneous P-modification as in (\ref{cmod.ksb.thm}.3).

Let $\bar{\mathbf D}$ be the birational transform of ${\mathbf D}$.
Since $\pi$ is small,
$$
K_{\bar {\mathbf S}}+d\bar{\mathbf D}\simq \pi^*\bigl(K_{{\mathbf S}}+d{\mathbf D}\bigr),
$$
and $(\bar {\mathbf S},d\bar{\mathbf D}+\bar {\mathbf S}_0)$ is also lc.
By  adjunction as in \cite[4.9]{kk-singbook}, we get that the pair
$\bigl(\bar {\mathbf S}_0, d\bar{\mathbf D}|_{\bar {\mathbf S}_0}\bigr)$ is also lc. Each divisor in $ d\bar{\mathbf D}|_{\bar {\mathbf S}_0}$ appears with a coefficient that is an integer multiple of $d$, and also at most $1$.
If $d>\frac12 $ then all coefficients are $d$ or $0$. If  $(S, dD)$ is klt
then all coefficients are $<1$, so again  they are  $d$ or $0$.

Since $K_{\bar {\mathbf S}_0}$ is $\pi_0$-ample,
$-\bar{\mathbf D}|_{\bar {\mathbf S}_0}$ is also $\pi_0$-ample,
hence its support must contain the whole exceptional divisor. Therefore
$\bar{\mathbf D}|_{\bar {\mathbf S}_0}$ is the sum of $\bar D$ plus the
$\pi_0$- exceptional divisor $\bar E_0$, so
$$
K_{\bar {\mathbf S}_0}+d(\bar D+\bar E_0)=\bigl(K_{\bar {\mathbf S}}+d\bar{\mathbf D}\bigr)|_{\bar {\mathbf S}_0}\simq \pi^*\bigl(K_{{\mathbf S}}+d{\mathbf D}\bigr).
 \qed
$$



 This takes us to the following, which is somewhat more general than needed for
 (\ref{K+D.KSB.P.thm}.2).


\begin{task} \label{task.task}
  Let $(S, dD)$ be an lc pair. Find all
  Q-modifications  $\pi:S'\to S$ with reduced exceptional divisor $E'$ such that
\begin{enumerate}
\item $K_{S'}$ is $\pi$-ample, and
\item $K_{S'}+d(D'+E')\simq 0$.
\end{enumerate}

Note that the conditions
$$
K_{S'}+dD'+\Delta\simq 0\qtq{and} \supp\Delta\subset E'
\eqno{(\ref{task.task}.3)}
$$
uniquely determine $\Delta$. If  (\ref{task.task}.2) holds then
$\Delta=dE'$, and each irreducible component of $E'$ gives a linear equation for $d$. We should thus expect no solutions  if there are several irreducible components.

We start by writing down solutions of (\ref{task.task}.3), and the using this to prove that there very few other solutions.
\end{task}




\begin{say}[Discrepancy divisors] \label{P.mod.P.mod.c.lem}
  There are some easy  solutions of (\ref{task.task}.3).

  \ref{P.mod.P.mod.c.lem}.1: $(S_{n,q}, D_{n,q})$.    Let $\pi:S'\to S_{n,q}$ be a $Q$-modification
  with reduced exceptional divisor $E'$. Then
  $K_{S'}+D'+E'\simq 0$.

  Indeed, this holds for the minimal resolution by (\ref{ord.of.K+C.say}.4), and continues to hold  since we blow up only nodes.

  \ref{P.mod.P.mod.c.lem}.2:  $(S^d_{n,q}, D^d_{n,q})$.
  On the minimal resolution $S^{\rm m}\to S:=S^d_{n,q}$, we have
  $$
  K_{S^{\rm m}}+D^{\rm m}_{n,q}+ \tfrac12(C'_0+C''_0)+\tsum _{i=1}^r C_i\simq 0
  $$
  by (\ref{ord.of.K+C.say}.5).
  As we blow up nodes, new curves appear with  half integer coefficients that are
  $\leq 1$.

  \ref{P.mod.P.mod.c.lem}.3: $(S_{n,q}, B_{n,q})$.  The formula is quite complicated already for the minimal resolution, see \cite[3.32--34]{kk-singbook}.
  We will not use these. Instead, we will choose a divisor  $\bar B_{n,q}$ so that $B_{n,q}+ \bar B_{n,q}=D_{n,q}$ and work with (\ref{P.mod.P.mod.c.lem}.1) instead.
\end{say}



The solution of  (\ref{task.task}) is especially simple for the pairs  $(S_{n,q}, dD_{n,q})$.

\begin{lem} \label{Dnq.triv.lem} Let $\pi:S'\to S:=S_{n,q}$ be a Q-modification
  with reduced exceptional curve $C'$ such that
  $K_{S'}$ is $\pi$-ample. Then
  \begin{enumerate}
  \item $K_{S'}+D'_{n,q}+C'\simq 0$, and
  \item if  $K_{S'}+d(D'+C')\simq 0$ for some $d<1$, then $\pi$ is an isomorphism.
    \end{enumerate}
\end{lem}

  Proof. The first claim was already noted in  (\ref{P.mod.P.mod.c.lem}.1).
  If $\pi$ has  $r\geq 1$ exceptional curve, then the  dual graph of $(S', D'+C')$ is
$$
   \bullet \ {-\boxed{\tfrac{n_0}{q_0}}-} \  c'_1 \  {-\boxed{\tfrac{n_1}{q_1}}-} \cdots  {-\boxed{\tfrac{n_{r-1}}{q_{r-1}}}-} \  c'_r \ {-\boxed{\tfrac{n_r}{q_r}}-} \ \bullet
   $$
   Since  $K_{S'}+D'+C'\simq 0$ and $K_{S'}+d(D'+C')\simq 0$, we get
   that $D'+C'\simq 0$ for  $d<1$. Then 
   $(D'+C')\cdot C'_1=0$ and  $K_{S'}\cdot C'_1=0$, a contradiction.\qed


\begin{cor} \label{34.1.no.defs.thm} A pair 
  $(S_{N,Q}, dD_{N,Q})$ is KSB smoothable iff $S_{N,Q}$ is KSB smoothable, hence
  $N=rn^2, Q=arn-1$. If these hold then
  the universal KSB deformation is given by
  $$
  \bigl((xy=z^{rn}+\tsum_{i=0}^{r-1}t_iz^{in}), d(z=0)\bigr)/\tfrac1{n}(1, -1, a, {\mathbf 0}).
  \hfill\qed
  %%\eqno{(\ref{34.1.no.defs.thm}.1)}
  $$
\end{cor}


Similarly, the dihedral pairs $(S^d_{n,q}, dD^d_{n,q})$ have no KSB smoothings for $d<1$.  The $d=1$ case is treated in (\ref{DP.char.lem}). 

    \begin{lem} \label{Cdnq.triv.lem} Let $\pi:S'\to S:=S^d_{n,q}$ be a Q-modification such that
     $K_{S'}$ is $\pi$-ample. Assume that $K_{S'}+d(D'+C')\simq 0$ for some
     $\frac12< d<1$. Then $\pi$ is an isomorphism.
\end{lem}

     Proof.   By (\ref{P.mod.P.mod.c.lem}.2) there are half integers
     $c_i\leq 1$ such that  $K_{S'}+D'+\sum_i c_iC'_i\simq 0$.
     Subtracting gives that
     $$
     (1-d)D'+\tsum (c_i-d)C'_i\simq 0.
     $$
     Since  $D'$ is $\pi$-nef, $-\tsum (c_i-d)C'_i$ is $\pi$-nef,
     hence $c_i-d\geq 0$ for every $i$.   Thus in fact $c_i=1$ for every $i$.
     The rest is now the same as in (\ref{Dnq.triv.lem}), but $S^d_{n,q}$ itself has no KSB smoothings. \qed






\medskip
By contrast, the pairs  $(S_{n,q}, dB_{n,q})$ have a much more complicated behavior.
   
   \begin{lem} \label{Bnq.almosttriv.lem} Let $\pi:S'\to S:=S_{n,q}$ be a Q-modification such that
     $K_{S'}$ is $\pi$-ample. Assume that $K_{S'}+d(B'+C')\simq 0$ for some
     $ d<1$. Then $\pi$ has at most 1 exceptional curve.
\end{lem}

   Proof.  If there are at least 2 exceptional curves, then 
   the dual graph  is
$$
   \bullet \ {-\boxed{\tfrac{n_0}{q_0}}-} \  c'_1 \  {-\boxed{\tfrac{n_1}{q_1}}-} c'_2 \ - \cdots
   $$
   By (\ref{P.mod.P.mod.c.lem}.3)  $K_{S'}+B'+C'+\bar B'\simq 0$.
   Subtracting, we get that
   $(1-d)(B'+C')+\bar B'\simq 0$.
   Here $\bar B'$ is disjoint from $C'_1$, so
   $(B'+C')\cdot C'_1=0$.  This in turn gives that 
$K_{S'}\cdot C'_1=0$, a contradiction.\qed


   
 
     
   


    
 


\begin{exmps} \label{4.0.smooth.exmp}
  Note that $S'$ with dual graph %%(\ref{1.excurve.lem}.1)
$$
   \bullet \ {-\boxed{\tfrac{n_0}{q_0}}-} \  c'_1 \  {-\boxed{\tfrac{n_1}{q_1}}}
   $$
  has a KSB deformation with Du~Val general fiber iff $S_{n_0, q_0}$ is a T-singularity and
   $S_{n_1, q_1}$ is a  Du~Val singularity. 
The simplest case  is the dual graph
     $$
   \bullet\ {-\boxed{\tfrac{4}{1}}-} \  c. 
    $$
   This gives that
   $
   \bigl(\a^2_{xy},\tfrac{2c-3}{2c-1}(y=0)\bigr) /\tfrac{1}{4c-1}(1, c)
     $
   is KSB smoothable.

   I found only 2 series whith KSB deformations for different values of $d$. The first is
   $ \bullet\ - \ 4\ - \  3  \ {-\boxed{A_{n+2}}}$
with P-modifications 
   $$
   \bullet\ {-\boxed{\tfrac{4}{1}}-} \  3  \ {-\boxed{A_{n+2}}}
   \qtq{and} 
   \bullet\ {-\boxed{\tfrac{18}{5}}-} \  2  \ {-\boxed{A_{n}}}.
   $$
  The second is  $ \bullet\ - \ 2\ - \ 5\ - \  3  \ {-\boxed{A_{n+1}}}$
with P-modifications 
   $$
   \bullet\ {-\boxed{\tfrac{9}{5}}-} \  3  \ {-\boxed{A_{n+1}}}
   \qtq{and} 
   \bullet\ {-\boxed{\tfrac{25}{14}}-} \  2  \ {-\boxed{A_{n}}}.
   $$
   \end{exmps}



\begin{exmp} \label{D.term.QC.exmp.1} Doubly KSB deformations of $D_n$ singularities are described as follows.
  For $r\geq 2$ the $D_r$-type
  pair
  $\bigl((z^2=x(y^2-x^{r-2})), (x=z=0)\bigr)$ has an 
$(r{-}2)$-parameter deformation 
  $$
  z^2=x\bigl(y^2-x^{r-2}-\tsum_{i=0}^{r-3} t_i x^i\bigr).
  $$
  (For $r=2$ we mean $(z^2=x(y^2-1))$, which has 2 singular points of type
  $\bigl((z^2=xy), (x=z=0)\bigr)$ at $(0, \pm 1, 0)$.)
  
  The divisor  $D:=(x=z=0)$ is $\q$-Cartier, since $2D=(x=0)$.
  So these are doubly KSB deformations.
The generic deformation  has 2 singular points of type
  $\bigl((z^2=xy), (x=z=0)\bigr)$ at   $(0,\pm \sqrt{-t_0}, 0)$.
  By (\ref{ord.of.K+C.say}.1) the latter has no KSB smoothings.
  \end{exmp}




\begin{exmp} \label{1/2.2.exmp} There are a several exceptional  KSB smoothable
  examples with $d=\frac12$. These are given by the dual graph, where $\ast$ can be either $\bullet$ or $2$
    $$
  \begin{array}{ccccc}
    && \ast &&\\
&& | &&\\
    \ast & - & 4  & - &   \ast\\
    && | &&\\
     && \ast &&
  \end{array}
  \qtq{where $D^{\rm m}+\Delta^{\rm m}$ is}
  \begin{array}{ccccc}
    && \tfrac12 &&\\
&& | &&\\
     \tfrac12 & - & 1  & - &    \tfrac12\\
    && | &&\\
     &&  \tfrac12 &&
  \end{array}
  $$
  Contracting the curve $(4)$ we get 
   doubly KSB smoothable examples.
\end{exmp}






\section{First order deformations}\label{first.ord.sec}


Various aspects of the first order deformation theory of cyclic quotients are worked out in \cite{k-alt}; the references below are to the version in
\cite[Sec.6.6]{k-modbook}.


Fix a pair  $(S, D):=\bigl(\a^2_{xy}, (xy=0)\bigr)/\tfrac1{n}(1,q)$.
Using  $\dx:=x(\partial/\partial x)$ and  $\dy:=y(\partial/\partial y)$,
the first order deformation space  $T^1(S)$ of $S$ has a basis consisting of certain operators
$$
D=\tfrac{a\dx-b\dy}{M},\qtq{where} a,b\in \c\qtq{and}  M=x^ry^s.
$$
Since  $\omega_S(D)\cong \o_S$, we can think of the curve $D$ as determined by
$$
\o_S\to \omega_S^{-1}\qtq{given by}  1\mapsto \tfrac{xy}{dx\wedge dy}.
$$
The lifting criterion \cite[6.58]{k-modbook} says that
$\frac{xy}{dx\wedge dy} $ lifts to a first order deformation determined by
 $D=(a\dx-b\dy)/(x^ry^s)$ iff
$$
D(xy)-xy\nabla D= (bs-ar) \tfrac{xy}{x^ry^s}
$$
equals $0$, as an element in $H^1(\a^2\setminus\{(0,0)\}, \o)$.
The latter holds if either $bs-ar=0$ or $\min\{r,s\}\leq 1$.

We can now look at the basis of  $T^1(S)$ given in \cite[6.73]{k-modbook}.
The above condition is satisfied by the basis elements $\dx/M_1, \dy/M_{r-1}$,
by a 1-dimensional subspace  of each  $\bigl\langle \dx/M_i, \dy/M_i\bigr\rangle$ for $2\leq i\leq r{-}2$, and by all the basis elements in  the list  \cite[6.73.2]{k-modbook}. Since the $r$ in \cite[6.73]{k-modbook} is $m(S)-1$ (where $m(S)$ is the multiplicity), we have proved the following.


\begin{thm}\label{KSB.codin.infini.thm}
  Set $(S, D):=\bigl(\a^2_{xy}, (xy=0)\bigr)/\tfrac1{n}(1,q)$. Then $T^1_{\rm KSB}(S, D)$ has codimension $m(S)-3$ in
  $T^1(S)$. \qed
  \end{thm}




\begin{exmp} We compute in more detail what happens for the singularity
  $S_{n,1}:=\a^2/\frac1{n}(1,1)$ for $n\geq 3$.

  The end result is that $\dim T^1(S_{n,1})=2n-4$, the Artin component is smooth
  of dimension $n-1$, and $T^1_{\rm KSB}(S_{n,1}, D_{n,1})$  also
has dimension $n-1$. Together, $T^1_{\rm Art}(S_{n,1})$ and $T^1_{\rm KSB}(S_{n,1}, D_{n,1})$
span $T^1(S_{n,1})$, so their intersection has dimension 2.

By Theorem~\ref{main.cyclic.thm}, $\defsp_{\rm KSB}(S_{n,1}, D_{n,1})$ has dimension 2; it is explicitly worked out  in   \cite[2.32]{k-modbook}.

To see these, note first  that by
\cite{pinkham}, the Artin component can be given by the  equations 
$$
\rank
\left(
\begin{array}{ccccc}
  z_0 & z_1 & \cdots &  z_{n-2} & z_{n-1}\\ 
 z_1-t_1 & z_2-t_2 & \cdots &  z_{n-1}-t_{n-1} & z_{n}
\end{array}
\right) \leq 1.
$$
The coordinate ring of $S_{n,1}$ is generated by the monomials
$M_i:=x^{n-i}y^i$. We claim that $T^1_{\rm Art}(S_n)$ is spanned
 by the vector fields
$$
D_i:=\tfrac{(n-i)\dx+i\dy}{nM_i}.
$$
We compute that 
$D_i(M_j)=M_j/M_i$,
which is invertible on the $x\neq 0$ chart for $j\geq i$. Thus, in this chart, we can choose 
$$
M_0, M_1, \dots, M_{i-1}, M_i+\epsilon Q_i,
M_{i+1}+\epsilon Q_{i+1},\dots,
M_n+\epsilon Q_n
$$
as liftings of the $M_i$, where we set  $Q_j:={M_j}/{M_i}$.
Note that  the matrix
$$
\left(
\begin{array}{ccccccccc}
  M_0 & M_1 & \cdots & M_{i-1} & M_i+\epsilon Q_i & \cdots & M_{n-2}+\epsilon Q_{n-2} & M_{n-1}+\epsilon Q_{n-1}\\
  M_1 & M_2 & \cdots & M_i & M_{i+1}+\epsilon Q_{i+1}&\cdots &  M_{n-1}+\epsilon Q_{n-1} & M_n+\epsilon Q_n
\end{array}
\right)
$$
has rank $\leq 1$. Relabeling the entries of the top row as
$z_0,\dots, z_{n-1}$, these become  the equations
$$
\rank
\left(
\begin{array}{ccccccccc}
  z_0 & z_1 & \cdots & z_{i-1} & z_i & \cdots &  z_{n-2} & z_{n-1}\\ 
 z_1 & z_2 & \cdots & z_{i}-\epsilon & z_{i+1} & \cdots &  z_{n-1} & z_{n}
\end{array}
\right) \leq 1.
$$
Thus these deformations span the Artin component.
\end{exmp}

The philosophy of \cite{k-modbook} says that in order to get the optimal moduli theory, we should  establish properties of 1-parameter deformations, and then formulate versions of these   properties  over arbitrary bases.
The second step is not necessarily unique. 
For deformations of pairs, this poses the following.


\begin{ques} Is there a good way to impose additional requirements for KSB deformations of a pair to make  $\defsp_{\rm KSB}(S_{n,q}, D_{n,q})$ smooth?
    \end{ques}



\section{Deformations of modifications}\label{def.mod.sect}


We prove that Q-modifications $X\to S$ have no (small or large)  KSB deformations over reduced bases that keep $S$ fixed. (See (\ref{An.mines.exmp}) for other deformations and non-reduced bases.)
More generally, we have the following.


\begin{thm} \label{rigid.quot.mod.thm}
  Let $S$ be a normal surface, $B$ a smooth, irreducible curve, and $g:Y\to S\times B$ a proper birational morphism, $Y$ normal. Assume that
  \begin{enumerate}
  \item $g_b:Y_b\to S$ is birational for every $b\in B$, and
    \item there is a dense,   open subset  $B^\circ\subset B$ such that, for every $b\in B^\circ$,
  $Y_b\to S$  is a Q-modification.
      \end{enumerate}
  Then, $Y\to B$ is trivial. That is, for any $b\in B$, 
    $$
  \bigl(Y\to S\times B\to B\bigr)\cong
  \bigl(Y_b\times B\to S\times B\to B\bigr).
  \eqno{(\ref{rigid.quot.mod.thm}.3)}
  $$
\end{thm}



Proof.   Over a dense,  open subset $B^*\subset B^\circ$,
the minimal resolutions   $\{Y_b^{\rm m}\to Y_b:b\in  B^*\}$
form a flat family. By  (\ref{rigid.quot.mod.lem.1}),
we have only finitely many choices for the centers of the blow-ups giving
$Y_b^{\rm m}\to S^{\rm m}$, so the $\{Y_b^{\rm m}:b\in  B^*\}$,
and hence also $\{ Y_b:b\in  B^*\}$, form localy trivial families.


Thus the family is also trivial over $B$
by (\ref{rigid.quot.mod.lem.2}). \qed




\begin{lem} \label{rigid.quot.mod.lem.2}
  Let $g:Y\to X\times B$ be a projective morphism such that
  $Y$ is normal and 
  $g_b:\red (Y_b)\to X_b$ is birational for every $b$.
  Assume that there is a dense, open $B^\circ\subset B$ such that
  $(Y^\circ\to B^\circ)\cong ((Y'\times B^\circ)\to B^\circ)$ for some
  normal $Y'$. Then $(Y\to B)\cong ((Y'\times B)\to B)$.
\end{lem}

Proof. Let $H$ be relatively ample on $Y\to X\times B$ and let $H'$ be its
birational  transform on $(Y'\times B)\to B$.
$H'$ is Cartier over $B^\circ$, so Cartier by \cite{ram-sam, sam-ram}; see also
\cite[4.21]{k-modbook}.
Thus $H'$ is relatively ample on $(Y'\times B)\to B$.
Thus $(Y\to B)\cong ((Y'\times B)\to B)$
 by   \cite{mats-mumf}; see \cite[11.39]{k-modbook} for the form that we use.  \qed 
 

 \begin{exmp}\label{An.mines.exmp} M-modifications have nontrivial deformations, as
   the next examples show.

   (\ref{An.mines.exmp}.1)  Let $S_P\to S$ be any P-modification whose construction involves blowing up a point  $p\in S^{\rm m}$. We can move the point
   $p$ along the exceptional curve and contract in the family to get a
   nontrivial flat deformation of $S_P\to S$ (keeping $S$ fixed).
   Here the canonical class of the generic fiber is not relatively nef.
   The simplest example is  the singularity $3\ -\ 3$. After one blow-up, the central and generic fibers are
   $$
   4\ -\ 1\ -\ 4 \qtq{resp.}   3\ -\ 4\ -\ 1.
   $$
   All curves marked $4$ are then contracted. The central fiber is an 
   M-modification. In the generic fiber, the canonical class has degree $-\frac12$ on the image of the
   curve marked $1$. 
   
    (\ref{An.mines.exmp}.2) For $A_n:=(xy=z^{n+1})$, the miniversal deformation of $A_n$ is
   $$
   \bigl(xy=z^{n+1}+\tsum_{i=0}^{n-1}t_iz^i\bigr),
   $$
   while the  miniversal deformation  of its minimal  resolution  is  obtained from 
   $$
   \bigl(xy=\tprod_{i=0}^{n}(z-s_i)\bigr), \qtq{(subject to $\tsum s_i=0$)}
   $$
   by  repeatedly blowing up the divisors $(x=z-s_i=0)$ as in (\ref{ksb.defs.say}.3).
   
   Let $\sigma_j$ be the $j$th elementary symmetric polynomial in the $s_i$.
   Then, over the Artin ring  $k[s_0,\dots, s_n]/(\sigma_0, \dots, \sigma_n)$, we get
   a   deformation  of the minimal  resolution which contracts to the
   trivial deformation of  $A_n$.

   A similar result holds for all  Du~Val singularities by \cite{art-bri},
   and for M-resolutions by \cite{be-ch}.
   \end{exmp}



 

 %%\bibliography{../refs-main/refs}

 \def\cprime{$'$} \def\cprime{$'$} \def\cprime{$'$} \def\cprime{$'$}
  \def\cprime{$'$} \def\dbar{\leavevmode\hbox to 0pt{\hskip.2ex
  \accent"16\hss}d} \def\cprime{$'$} \def\cprime{$'$}
  \def\polhk#1{\setbox0=\hbox{#1}{\ooalign{\hidewidth
  \lower1.5ex\hbox{`}\hidewidth\crcr\unhbox0}}} \def\cprime{$'$}
  \def\cprime{$'$} \def\cprime{$'$} \def\cprime{$'$}
  \def\polhk#1{\setbox0=\hbox{#1}{\ooalign{\hidewidth
  \lower1.5ex\hbox{`}\hidewidth\crcr\unhbox0}}} \def\cdprime{$''$}
  \def\cprime{$'$} \def\cprime{$'$} \def\cprime{$'$} \def\cprime{$'$}
\providecommand{\bysame}{\leavevmode\hbox to3em{\hrulefill}\thinspace}
\providecommand{\MR}{\relax\ifhmode\unskip\space\fi MR }
% \MRhref is called by the amsart/book/proc definition of \MR.
\providecommand{\MRhref}[2]{%
  \href{http://www.ams.org/mathscinet-getitem?mr=#1}{#2}
}
\providecommand{\href}[2]{#2}
\begin{thebibliography}{DRU22}

\bibitem[AK19]{k-alt}
Klaus Altmann and J{\'a}nos Koll{\'a}r, \emph{The dualizing sheaf on
  first-order deformations of toric surface singularities}, J. Reine Angew.
  Math. \textbf{753} (2019), 137--158. \MR{3987866}

\bibitem[Art74]{art-bri}
Michael Artin, \emph{Algebraic construction of {B}rieskorn's resolutions}, J.
  Algebra \textbf{29} (1974), 330--348.

\bibitem[BC94]{be-ch}
Kurt Behnke and Jan~Arthur Christophersen, \emph{{$M$}-resolutions and
  deformations of quotient singularities}, Amer. J. Math. \textbf{116} (1994),
  no.~4, 881--903. \MR{1287942}

\bibitem[Bin87]{bing}
J\"urgen Bingener, \emph{Lokake {M}odulr\"aume in den analytischen {G}eometrie
  {I}--{II}}, Aspects of math., vol. 302, Viehweg, Braunschweig, 1987.

\bibitem[BPV84]{bpv}
W.~Barth, C.~Peters, and A.~V{an de Ven}, \emph{Compact complex surfaces},
  Ergebnisse der Mathematik und ihrer Grenzgebiete (3), vol.~4,
  Springer-Verlag, Berlin, 1984.

\bibitem[Bri68]{briesk}
Egbert Brieskorn, \emph{Die {A}ufl\"osung der rationalen {S}ingularit\"aten
  holomorpher {A}bbildungen}, Math. Ann. \textbf{178} (1968), 255--270.
  \MR{0233819 (38 \#2140)}

\bibitem[BS11]{MR2843099}
Mohan Bhupal and Andr\'{a}s~I. Stipsicz, \emph{Weighted homogeneous
  singularities and rational homology disk smoothings}, Amer. J. Math.
  \textbf{133} (2011), no.~5, 1259--1297. \MR{2843099}

\bibitem[dJvS91]{MR1135880}
Theo de~Jong and Duco van Straten, \emph{On the base space of a semi-universal
  deformation of rational quadruple points}, Ann. of Math. (2) \textbf{134}
  (1991), no.~3, 653--678. \MR{1135880}

\bibitem[DRU22]{MR4502069}
Eduardo Dias, Carlos Rito, and Giancarlo Urz\'{u}a, \emph{On degenerations of
  {$\mathbb{Z}/2$}-{G}odeaux surfaces}, Rev. Mat. Iberoam. \textbf{38} (2022),
  no.~5, 1399--1423. \MR{4502069}

\bibitem[HM94]{MR1320605}
Herwig Hauser and Gerd M{\"u}ller, \emph{A rank theorem for analytic maps
  between power series spaces}, Inst. Hautes \'Etudes Sci. Publ. Math. (1994),
  no.~80, 95--115 (1995). \MR{1320605 (96b:46065)}

\bibitem[JS23]{jeon2023deformations}
Jaekwan Jeon and Dongsoo Shin, \emph{Deformations of weighted homogeneous
  surface singularities with big central node}, 2023.

\bibitem[Jun08]{jung}
Heinrich W.~E. Jung, \emph{Darstellung der {F}unktionen eines algebraischen
  {K}\"orpers zweier unabh\"angigen {V}er\"anderlichen {$x,y$} in der
  {U}mgebung einer {S}telle {$x=a,\ y=b$}}, J. Reine Angew. Math. \textbf{133}
  (1908), 289--314. \MR{1580742}

\bibitem[KK23]{k-db3}
J{\'a}nos Koll{\'a}r and S{\'a}ndor~J Kov{\'a}cs, \emph{{KSB} stability is
  automatic in codimension 3}, 2023, arXiv:2304.09009.

\bibitem[KLP12]{MR2995167}
JongHae Keum, Yongnam Lee, and Heesang Park, \emph{Construction of surfaces of
  general type from elliptic surfaces via {$\mathbb{Q}$}-{G}orenstein
  smoothing}, Math. Z. \textbf{272} (2012), no.~3-4, 1243--1257. \MR{2995167}

\bibitem[KM92]{km-flips}
J{\'a}nos Koll{\'a}r and Shigefumi Mori, \emph{Classification of
  three-dimensional flips}, J. Amer. Math. Soc. \textbf{5} (1992), no.~3,
  533--703. \MR{1149195 (93i:14015)}

\bibitem[Kol91]{k-etc}
J{\'a}nos Koll{\'a}r, \emph{Flips, flops, minimal models, etc}, Surveys in
  differential geometry ({C}ambridge, {MA}, 1990), Lehigh Univ., Bethlehem, PA,
  1991, pp.~113--199. \MR{1144527}

\bibitem[Kol07]{k-res}
\bysame, \emph{Lectures on resolution of singularities}, Annals of Mathematics
  Studies, vol. 166, Princeton University Press, Princeton, NJ, 2007.

\bibitem[Kol13]{kk-singbook}
\bysame, \emph{Singularities of the minimal model program}, Cambridge Tracts in
  Mathematics, vol. 200, Cambridge University Press, Cambridge, 2013, With the
  collaboration of S{\'a}ndor Kov{\'a}cs.

\bibitem[Kol23]{k-modbook}
\bysame, \emph{Families of varieties of general type}, Cambridge Tracts in
  Mathematics, vol. 231, Cambridge University Press, Cambridge, 2023, With the
  collaboration of Klaus Altmann and S{\'a}ndor Kov{\'a}cs. \MR{4566297}

\bibitem[KSB88]{ksb}
J{\'a}nos Koll{\'a}r and N.~I. Shepherd-Barron, \emph{Threefolds and
  deformations of surface singularities}, Invent. Math. \textbf{91} (1988),
  no.~2, 299--338.

\bibitem[Lau73a]{MR367277}
Henry~B. Laufer, \emph{Deformations of resolutions of two-dimensional
  singularities}, Rice Univ. Stud. \textbf{59} (1973), no.~1, 53--96.
  \MR{367277}

\bibitem[Lau73b]{MR333238}
\bysame, \emph{Taut two-dimensional singularities}, Math. Ann. \textbf{205}
  (1973), 131--164. \MR{333238}

\bibitem[LP07]{lee-park}
Yongnam Lee and Jongil Park, \emph{A simply connected surface of general type
  with {$p_g=0$} and {$K^2=2$}}, Invent. Math. \textbf{170} (2007), no.~3,
  483--505. \MR{2357500 (2008m:14076)}

\bibitem[MM64]{mats-mumf}
Teruhisa Matsusaka and David Mumford, \emph{Two fundamental theorems on
  deformations of polarized varieties}, Amer. J. Math. \textbf{86} (1964),
  668--684. \MR{0171778 (30 \#2005)}

\bibitem[Pin77]{pinkham}
H.~Pinkham, \emph{Normal surface singularities with {$C\sp*$} action}, Math.
  Ann. \textbf{227} (1977), no.~2, 183--193. \MR{55 \#5623}

\bibitem[PPS23]{pa-pa-sh-2}
Heesang Park, Jongil Park, and Dongsoo Shin, \emph{Surfaces of general type
  with $p_g=1$ and $q=0$}, J. Korean Math. Soc. \textbf{50} (2023), no.~3,
  493--507.

\bibitem[PS22]{park2022deformations}
Heesang Park and Dongsoo Shin, \emph{Deformations of sandwiched surface
  singularities and the minimal model program}, 2022.

\bibitem[Ram63]{ram-sam}
Chakravarthi~Padmanabhan Ramanujam, \emph{Appendix to: {Q}uotient space by an
  abelian variety, by {C}.{S}.~{S}eshadri}, Math. Ann. \textbf{152} (1963),
  185--194.

\bibitem[Rei93]{reid-cyclic}
Miles Reid, \emph{Surface cyclic quotient singularities and {H}irzebruch-{J}ung
  resolutions}, 1993, Available at
  \url{https://homepages.warwick.ac.uk/~masda/surf/more/cyclic.pdf}.

\bibitem[RU22]{MR4509033}
Javier Reyes and Giancarlo Urz\'{u}a, \emph{Rational configurations in {K}3
  surfaces and simply-connected {$p_g=1$} surfaces with
  {$K^2=1,2,3,4,5,6,7,8,9$}}, Math. Z. \textbf{302} (2022), no.~4, 2435--2467.
  \MR{4509033}

\bibitem[Sam62]{sam-ram}
Pierre Samuel, \emph{Sur une conjecture de {G}rothendieck}, C. R. Acad. Sci.
  Paris S\'{e}r. A-B \textbf{255} (1962), 3101--3103.

\bibitem[Ser06]{sernesi}
Edoardo Sernesi, \emph{Deformations of algebraic schemes}, Grundlehren der
  Mathematischen Wissenschaften, vol. 334, Springer-Verlag, Berlin, 2006.
  \MR{2247603 (2008e:14011)}

\bibitem[Ste91a]{MR1129040}
Jan Stevens, \emph{On the versal deformation of cyclic quotient singularities},
  Singularity theory and its applications, {P}art {I} ({C}oventry, 1988/1989),
  Lecture Notes in Math., vol. 1462, Springer, Berlin, 1991, pp.~302--319.
  \MR{1129040 (93b:14061)}

\bibitem[Ste91b]{MR1094706}
\bysame, \emph{Partial resolutions of rational quadruple points}, Internat. J.
  Math. \textbf{2} (1991), no.~2, 205--221. \MR{1094706}

\bibitem[Ste93]{MR1213355}
\bysame, \emph{Partial resolutions of quotient singularities}, Manuscripta
  Math. \textbf{79} (1993), no.~1, 7--11. \MR{1213355}

\bibitem[Wah76]{MR54:10261}
Jonathan Wahl, \emph{Equisingular deformations of normal surface singularities.
  {I}}, Ann. of Math. (2) \textbf{104} (1976), no.~2, 325--356. \MR{MR54
  \#10261}

\bibitem[Wah81]{MR83h:14029}
\bysame, \emph{Smoothings of normal surface singularities}, Topology
  \textbf{20} (1981), no.~3, 219--246. \MR{MR83h:14029}

\bibitem[Wah21]{MR4270358}
\bysame, \emph{Complex surface singularities with rational homology disk
  smoothings}, Singularities and their interaction with geometry and low
  dimensional topology---in honor of {A}ndr\'{a}s {N}\'{e}methi, Trends Math.,
  Birkh\"{a}user/Springer, Cham, [2021] \copyright 2021, pp.~259--287.
  \MR{4270358}

\end{thebibliography}


 \bigskip


  Princeton University, Princeton NJ 08544-1000, \

  \email{kollar@math.princeton.edu}



 \end{document}


